\section{Circuitos eléctricos y electrónica}

\subsection*{H-001 \quad Respuesta B}\label{s:H-001}
\textbf{Razonamiento.} Por Ohm, I=12/3000=0.004 A=4 mA.

\textbf{A) Distractor. }Invierte la relación entre resistencia y voltaje.

\textbf{C) Distractor. }Multiplicar voltaje y resistencia no produce corriente.

\textbf{Repaso:} \hyperref[sec:h1]{Circuitos eléctricos y electrónica}. \textbf{Fuentes:} \cite{openstax,electronics,digital}. \hyperref[q:H-001]{Volver al reactivo}.\par\medskip

\subsection*{H-002 \quad Respuesta C}\label{s:H-002}
\textbf{Razonamiento.} Req=(6×3)/(6+3)=2 kΩ.

\textbf{A) Distractor. }La suma directa corresponde a serie.

\textbf{B) Distractor. }El promedio aritmético no es la regla del paralelo.

\textbf{Repaso:} \hyperref[sec:h1]{Circuitos eléctricos y electrónica}. \textbf{Fuentes:} \cite{openstax,electronics,digital}. \hyperref[q:H-002]{Volver al reactivo}.\par\medskip

\subsection*{H-003 \quad Respuesta B}\label{s:H-003}
\textbf{Razonamiento.} La conservación de carga exige 8=3+I, de donde I=5 mA.

\textbf{A) Distractor. }El producto no satisface la ley de corrientes.

\textbf{C) Distractor. }Suma una corriente que ya sale del nodo.

\textbf{Repaso:} \hyperref[sec:h1]{Circuitos eléctricos y electrónica}. \textbf{Fuentes:} \cite{openstax,electronics,digital}. \hyperref[q:H-003]{Volver al reactivo}.\par\medskip

\subsection*{H-004 \quad Respuesta C}\label{s:H-004}
\textbf{Razonamiento.} El potencial eléctrico se mide en joules por coulomb, equivalentes a volts.

\textbf{A) Distractor. }La corriente es carga por unidad de tiempo.

\textbf{B) Distractor. }La potencia es energía por unidad de tiempo.

\textbf{Repaso:} \hyperref[sec:h1]{Circuitos eléctricos y electrónica}. \textbf{Fuentes:} \cite{openstax,electronics,digital}. \hyperref[q:H-004]{Volver al reactivo}.\par\medskip

\subsection*{H-005 \quad Respuesta A}\label{s:H-005}
\textbf{Razonamiento.} XOR vale uno para 01 y 10.

\textbf{B) Distractor. }El caso 11 corresponde a la condición de AND.

\textbf{C) Distractor. }El caso 00 hace cero a XOR.

\textbf{Repaso:} \hyperref[sec:h1]{Circuitos eléctricos y electrónica}. \textbf{Fuentes:} \cite{openstax,electronics,digital}. \hyperref[q:H-005]{Volver al reactivo}.\par\medskip

\subsection*{H-006 \quad Respuesta C}\label{s:H-006}
\textbf{Razonamiento.} El inductor almacena energía LI²/2.

\textbf{A) Distractor. }El capacitor almacena energía en un campo eléctrico.

\textbf{B) Distractor. }El resistor ideal disipa energía.

\textbf{Repaso:} \hyperref[sec:h1]{Circuitos eléctricos y electrónica}. \textbf{Fuentes:} \cite{openstax,electronics,digital}. \hyperref[q:H-006]{Volver al reactivo}.\par\medskip

\subsection*{H-007 \quad Respuesta A}\label{s:H-007}
\textbf{Razonamiento.} La memoria o estado introduce dependencia de la historia.

\textbf{B) Distractor. }Los circuitos secuenciales sí pueden construirse con compuertas.

\textbf{C) Distractor. }Describe al circuito combinacional.

\textbf{Repaso:} \hyperref[sec:h1]{Circuitos eléctricos y electrónica}. \textbf{Fuentes:} \cite{openstax,electronics,digital}. \hyperref[q:H-007]{Volver al reactivo}.\par\medskip

\subsection*{H-008 \quad Respuesta B}\label{s:H-008}
\textbf{Razonamiento.} Vsal=9×4/(2+4)=6 V.

\textbf{A) Distractor. }Existe caída en ambos resistores.

\textbf{C) Distractor. }Es la caída sobre el resistor superior.

\textbf{Repaso:} \hyperref[sec:h1]{Circuitos eléctricos y electrónica}. \textbf{Fuentes:} \cite{openstax,electronics,digital}. \hyperref[q:H-008]{Volver al reactivo}.\par\medskip

\subsection*{H-009 \quad Respuesta A}\label{s:H-009}
\textbf{Razonamiento.} P=I²R=0.04×100=4 W.

\textbf{B) Distractor. }IR entrega voltaje, no potencia.

\textbf{C) Distractor. }Falta un factor de diez respecto al cálculo correcto.

\textbf{Repaso:} \hyperref[sec:h1]{Circuitos eléctricos y electrónica}. \textbf{Fuentes:} \cite{openstax,electronics,digital}. \hyperref[q:H-009]{Volver al reactivo}.\par\medskip

\subsection*{H-010 \quad Respuesta B}\label{s:H-010}
\textbf{Razonamiento.} τ=RC=10000×0.00001=0.1 s.

\textbf{A) Distractor. }Introduce dos órdenes de magnitud de error.

\textbf{C) Distractor. }No convierte microfarads a farads correctamente.

\textbf{Repaso:} \hyperref[sec:h1]{Circuitos eléctricos y electrónica}. \textbf{Fuentes:} \cite{openstax,electronics,digital}. \hyperref[q:H-010]{Volver al reactivo}.\par\medskip

\subsection*{H-011 \quad Respuesta C}\label{s:H-011}
\textbf{Razonamiento.} En realimentación negativa y operación lineal, Av=1+Rf/Rg=4.

\textbf{A) Distractor. }Omite la unidad del modelo no inversor.

\textbf{B) Distractor. }Corresponde a otra configuración y signo.

\textbf{Repaso:} \hyperref[sec:h1]{Circuitos eléctricos y electrónica}. \textbf{Fuentes:} \cite{openstax,electronics,digital}. \hyperref[q:H-011]{Volver al reactivo}.\par\medskip

\subsection*{H-012 \quad Respuesta B}\label{s:H-012}
\textbf{Razonamiento.} El resistor absorbe 5−2=3 V; R=3/0.01=300 Ω.

\textbf{A) Distractor. }Utiliza la caída del LED como caída del resistor.

\textbf{C) Distractor. }Ignora la caída del LED.

\textbf{Repaso:} \hyperref[sec:h1]{Circuitos eléctricos y electrónica}. \textbf{Fuentes:} \cite{openstax,electronics,digital}. \hyperref[q:H-012]{Volver al reactivo}.\par\medskip

\subsection*{H-013 \quad Respuesta C}\label{s:H-013}
\textbf{Razonamiento.} La ley de De Morgan cambia AND por OR y niega ambas entradas.

\textbf{A) Distractor. }Mantiene incorrectamente la operación AND.

\textbf{B) Distractor. }Faltan las negaciones de las entradas.

\textbf{Repaso:} \hyperref[sec:h1]{Circuitos eléctricos y electrónica}. \textbf{Fuentes:} \cite{openstax,electronics,digital}. \hyperref[q:H-013]{Volver al reactivo}.\par\medskip

\subsection*{H-014 \quad Respuesta A}\label{s:H-014}
\textbf{Razonamiento.} Se necesita 2\textasciicircum{}n≥101; 64 no basta y 128 sí.

\textbf{B) Distractor. }Seis bits representan únicamente 0 a 63.

\textbf{C) Distractor. }No se necesita un bit por valor.

\textbf{Repaso:} \hyperref[sec:h1]{Circuitos eléctricos y electrónica}. \textbf{Fuentes:} \cite{openstax,electronics,digital}. \hyperref[q:H-014]{Volver al reactivo}.\par\medskip

\subsection*{H-015 \quad Respuesta C}\label{s:H-015}
\textbf{Razonamiento.} El flip-flop D copia D en el flanco activo.

\textbf{A) Distractor. }La inversión corresponde a comportamiento de conmutación, no de D.

\textbf{B) Distractor. }No describe la captura de un uno.

\textbf{Repaso:} \hyperref[sec:h1]{Circuitos eléctricos y electrónica}. \textbf{Fuentes:} \cite{openstax,electronics,digital}. \hyperref[q:H-015]{Volver al reactivo}.\par\medskip

\subsection*{H-016 \quad Respuesta A}\label{s:H-016}
\textbf{Razonamiento.} El modelo aproxima la región de conducción directa mediante una caída fija.

\textbf{B) Distractor. }Un diodo polarizado directamente no permanece siempre abierto.

\textbf{C) Distractor. }No es el modelo de caída constante.

\textbf{Repaso:} \hyperref[sec:h1]{Circuitos eléctricos y electrónica}. \textbf{Fuentes:} \cite{openstax,electronics,digital}. \hyperref[q:H-016]{Volver al reactivo}.\par\medskip

\subsection*{H-017 \quad Respuesta B}\label{s:H-017}
\textbf{Razonamiento.} Primero se fija tierra, luego se escriben ecuaciones de corriente y se resuelven.

\textbf{A) Distractor. }KCL requiere definir previamente las tensiones respecto a una referencia.

\textbf{C) Distractor. }No se resuelve antes de formular el sistema.

\textbf{Repaso:} \hyperref[sec:h1]{Circuitos eléctricos y electrónica}. \textbf{Fuentes:} \cite{openstax,electronics,digital}. \hyperref[q:H-017]{Volver al reactivo}.\par\medskip

\subsection*{H-018 \quad Respuesta A}\label{s:H-018}
\textbf{Razonamiento.} AND exige ambas, OR al menos una y NOT invierte.

\textbf{B) Distractor. }Asigna funciones incorrectas a las tres operaciones.

\textbf{C) Distractor. }Confunde negación y operaciones de dos entradas.

\textbf{Repaso:} \hyperref[sec:h1]{Circuitos eléctricos y electrónica}. \textbf{Fuentes:} \cite{openstax,electronics,digital}. \hyperref[q:H-018]{Volver al reactivo}.\par\medskip

\subsection*{H-019 \quad Respuesta B}\label{s:H-019}
\textbf{Razonamiento.} El equivalente de Thévenin forma un divisor de dos resistencias iguales.

\textbf{A) Distractor. }La carga no está cortocircuitada.

\textbf{C) Distractor. }Supone una fuente ideal sin resistencia interna.

\textbf{Repaso:} \hyperref[sec:h1]{Circuitos eléctricos y electrónica}. \textbf{Fuentes:} \cite{openstax,electronics,digital}. \hyperref[q:H-019]{Volver al reactivo}.\par\medskip

\subsection*{H-020 \quad Respuesta C}\label{s:H-020}
\textbf{Razonamiento.} La resistencia de polarización fija el nivel cuando el interruptor está abierto.

\textbf{A) Distractor. }La resolución analógica no fija un nodo flotante.

\textbf{B) Distractor. }Puede crear contención eléctrica y no resuelve la entrada.

\textbf{Repaso:} \hyperref[sec:h1]{Circuitos eléctricos y electrónica}. \textbf{Fuentes:} \cite{openstax,electronics,digital}. \hyperref[q:H-020]{Volver al reactivo}.\par\medskip

\subsection*{H-021 \quad Respuesta B}\label{s:H-021}
\textbf{Razonamiento.} La salida está limitada por la alimentación y el margen del dispositivo.

\textbf{A) Distractor. }El cálculo de ganancia solo vale en la región lineal.

\textbf{C) Distractor. }No puede generar arbitrariamente una tensión fuera de sus límites.

\textbf{Repaso:} \hyperref[sec:h1]{Circuitos eléctricos y electrónica}. \textbf{Fuentes:} \cite{openstax,electronics,digital}. \hyperref[q:H-021]{Volver al reactivo}.\par\medskip

\subsection*{H-022 \quad Respuesta C}\label{s:H-022}
\textbf{Razonamiento.} Diferencias de propagación pueden crear transitorios no previstos por la tabla estática.

\textbf{A) Distractor. }La retención sin alimentación no causa ese pulso.

\textbf{B) Distractor. }No hay un sistema mecánico en el caso.

\textbf{Repaso:} \hyperref[sec:h1]{Circuitos eléctricos y electrónica}. \textbf{Fuentes:} \cite{openstax,electronics,digital}. \hyperref[q:H-022]{Volver al reactivo}.\par\medskip

\section{Arquitectura de computadoras y su organización}

\subsection*{H-023 \quad Respuesta C}\label{s:H-023}
\textbf{Razonamiento.} El PC contiene la dirección que guía la siguiente búsqueda de instrucción.

\textbf{A) Distractor. }La capacidad de RAM es una propiedad distinta.

\textbf{B) Distractor. }Un registro arquitectónico no mide tensión.

\textbf{Repaso:} \hyperref[sec:h2]{Arquitectura de computadoras y su organización}. \textbf{Fuentes:} \cite{intel,architecture}. \hyperref[q:H-023]{Volver al reactivo}.\par\medskip

\subsection*{H-024 \quad Respuesta A}\label{s:H-024}
\textbf{Razonamiento.} 32+8+4+1=45.

\textbf{B) Distractor. }Omite el bit que aporta 4.

\textbf{C) Distractor. }Añade un peso que no está activo.

\textbf{Repaso:} \hyperref[sec:h2]{Arquitectura de computadoras y su organización}. \textbf{Fuentes:} \cite{intel,architecture}. \hyperref[q:H-024]{Volver al reactivo}.\par\medskip

\subsection*{H-025 \quad Respuesta C}\label{s:H-025}
\textbf{Razonamiento.} La inversión de 00000001 más uno produce 11111111.

\textbf{A) Distractor. }Ese patrón representa −127 en complemento a dos.

\textbf{B) Distractor. }Representa +1.

\textbf{Repaso:} \hyperref[sec:h2]{Arquitectura de computadoras y su organización}. \textbf{Fuentes:} \cite{intel,architecture}. \hyperref[q:H-025]{Volver al reactivo}.\par\medskip

\subsection*{H-026 \quad Respuesta A}\label{s:H-026}
\textbf{Razonamiento.} DRAM requiere alimentación y refresco para retener datos.

\textbf{B) Distractor. }La ROM conserva su contenido sin alimentación.

\textbf{C) Distractor. }Flash es memoria no volátil.

\textbf{Repaso:} \hyperref[sec:h2]{Arquitectura de computadoras y su organización}. \textbf{Fuentes:} \cite{intel,architecture}. \hyperref[q:H-026]{Volver al reactivo}.\par\medskip

\subsection*{H-027 \quad Respuesta B}\label{s:H-027}
\textbf{Razonamiento.} La unidad aritmético-lógica ejecuta estas operaciones.

\textbf{A) Distractor. }Proporciona energía, no ejecuta instrucciones.

\textbf{C) Distractor. }Su función principal es gestionar la visualización.

\textbf{Repaso:} \hyperref[sec:h2]{Arquitectura de computadoras y su organización}. \textbf{Fuentes:} \cite{intel,architecture}. \hyperref[q:H-027]{Volver al reactivo}.\par\medskip

\subsection*{H-028 \quad Respuesta A}\label{s:H-028}
\textbf{Razonamiento.} La localidad temporal describe accesos repetidos en un intervalo corto.

\textbf{B) Distractor. }El medio óptico no define este principio.

\textbf{C) Distractor. }La contigüidad corresponde a localidad espacial.

\textbf{Repaso:} \hyperref[sec:h2]{Arquitectura de computadoras y su organización}. \textbf{Fuentes:} \cite{intel,architecture}. \hyperref[q:H-028]{Volver al reactivo}.\par\medskip

\subsection*{H-029 \quad Respuesta B}\label{s:H-029}
\textbf{Razonamiento.} La arquitectura del conjunto de instrucciones es el contrato de programación de la CPU.

\textbf{A) Distractor. }El sistema operativo utiliza la ISA, pero no la sustituye.

\textbf{C) Distractor. }Diferentes microarquitecturas pueden implementar la misma ISA.

\textbf{Repaso:} \hyperref[sec:h2]{Arquitectura de computadoras y su organización}. \textbf{Fuentes:} \cite{intel,architecture}. \hyperref[q:H-029]{Volver al reactivo}.\par\medskip

\subsection*{H-030 \quad Respuesta C}\label{s:H-030}
\textbf{Razonamiento.} Hay 2\textasciicircum{}16 bytes=65536 bytes=64 KiB.

\textbf{A) Distractor. }Confunde bits de dirección con capacidad.

\textbf{B) Distractor. }Omite la exponenciación y la unidad byte.

\textbf{Repaso:} \hyperref[sec:h2]{Arquitectura de computadoras y su organización}. \textbf{Fuentes:} \cite{intel,architecture}. \hyperref[q:H-030]{Volver al reactivo}.\par\medskip

\subsection*{H-031 \quad Respuesta B}\label{s:H-031}
\textbf{Razonamiento.} Tiempo=IC×CPI/f=6×10\textasciicircum{}6/10\textasciicircum{}9=0.006 s.

\textbf{A) Distractor. }No convierte la frecuencia de GHz.

\textbf{C) Distractor. }Ignora los tres ciclos por instrucción.

\textbf{Repaso:} \hyperref[sec:h2]{Arquitectura de computadoras y su organización}. \textbf{Fuentes:} \cite{intel,architecture}. \hyperref[q:H-031]{Volver al reactivo}.\par\medskip

\subsection*{H-032 \quad Respuesta C}\label{s:H-032}
\textbf{Razonamiento.} AMAT=2+0.05×80=6 ns.

\textbf{A) Distractor. }Aplica la penalización a todos los accesos.

\textbf{B) Distractor. }Cuenta la penalización media pero omite el acceso inicial.

\textbf{Repaso:} \hyperref[sec:h2]{Arquitectura de computadoras y su organización}. \textbf{Fuentes:} \cite{intel,architecture}. \hyperref[q:H-032]{Volver al reactivo}.\par\medskip

\subsection*{H-033 \quad Respuesta A}\label{s:H-033}
\textbf{Razonamiento.} Little-endian coloca el byte menos significativo primero.

\textbf{B) Distractor. }Sería el primer byte en big-endian.

\textbf{C) Distractor. }Es el segundo byte desde el menos significativo.

\textbf{Repaso:} \hyperref[sec:h2]{Arquitectura de computadoras y su organización}. \textbf{Fuentes:} \cite{intel,architecture}. \hyperref[q:H-033]{Volver al reactivo}.\par\medskip

\subsection*{H-034 \quad Respuesta C}\label{s:H-034}
\textbf{Razonamiento.} 260 módulo 256 es 4.

\textbf{A) Distractor. }No cabe en 8 bits.

\textbf{B) Distractor. }−6 corresponde a interpretar 250 como un valor con signo antes de sumar.

\textbf{Repaso:} \hyperref[sec:h2]{Arquitectura de computadoras y su organización}. \textbf{Fuentes:} \cite{intel,architecture}. \hyperref[q:H-034]{Volver al reactivo}.\par\medskip

\subsection*{H-035 \quad Respuesta A}\label{s:H-035}
\textbf{Razonamiento.} DMA delega la transferencia a un controlador y coordina su finalización.

\textbf{B) Distractor. }El sondeo consume intervención repetida de la CPU.

\textbf{C) Distractor. }Sin memoria no se logra almacenar el bloque.

\textbf{Repaso:} \hyperref[sec:h2]{Arquitectura de computadoras y su organización}. \textbf{Fuentes:} \cite{intel,architecture}. \hyperref[q:H-035]{Volver al reactivo}.\par\medskip

\subsection*{H-036 \quad Respuesta B}\label{s:H-036}
\textbf{Razonamiento.} La dependencia entre productor y consumidor puede requerir forwarding o espera.

\textbf{A) Distractor. }La localidad se refiere a patrones de acceso.

\textbf{C) Distractor. }El conflicto surge de operandos, no de temperatura.

\textbf{Repaso:} \hyperref[sec:h2]{Arquitectura de computadoras y su organización}. \textbf{Fuentes:} \cite{intel,architecture}. \hyperref[q:H-036]{Volver al reactivo}.\par\medskip

\subsection*{H-037 \quad Respuesta A}\label{s:H-037}
\textbf{Razonamiento.} Se reservan 4 bits de desplazamiento y 4 de índice; quedan 8 de etiqueta.

\textbf{B) Distractor. }Cuatro son para cada uno de los otros campos.

\textbf{C) Distractor. }Omite descontar uno de los dos campos.

\textbf{Repaso:} \hyperref[sec:h2]{Arquitectura de computadoras y su organización}. \textbf{Fuentes:} \cite{intel,architecture}. \hyperref[q:H-037]{Volver al reactivo}.\par\medskip

\subsection*{H-038 \quad Respuesta B}\label{s:H-038}
\textbf{Razonamiento.} El retorno exige restaurar el flujo y el estado relevante según la arquitectura.

\textbf{A) Distractor. }No es necesario copiar todo el almacenamiento para cada interrupción.

\textbf{C) Distractor. }El nombre no permite reconstruir registros ni punto de ejecución.

\textbf{Repaso:} \hyperref[sec:h2]{Arquitectura de computadoras y su organización}. \textbf{Fuentes:} \cite{intel,architecture}. \hyperref[q:H-038]{Volver al reactivo}.\par\medskip

\subsection*{H-039 \quad Respuesta C}\label{s:H-039}
\textbf{Razonamiento.} La CPU obtiene la instrucción, interpreta sus campos y realiza la operación.

\textbf{A) Distractor. }No puede decodificar una instrucción aún no obtenida.

\textbf{B) Distractor. }Ejecutar requiere conocer antes la instrucción.

\textbf{Repaso:} \hyperref[sec:h2]{Arquitectura de computadoras y su organización}. \textbf{Fuentes:} \cite{intel,architecture}. \hyperref[q:H-039]{Volver al reactivo}.\par\medskip

\subsection*{H-040 \quad Respuesta B}\label{s:H-040}
\textbf{Razonamiento.} Los registros sostienen operandos, la caché reduce accesos lentos y el SSD persiste datos.

\textbf{A) Distractor. }Confunde niveles y funciones de la jerarquía.

\textbf{C) Distractor. }Asigna persistencia al registro y funciones incorrectas al resto.

\textbf{Repaso:} \hyperref[sec:h2]{Arquitectura de computadoras y su organización}. \textbf{Fuentes:} \cite{intel,architecture}. \hyperref[q:H-040]{Volver al reactivo}.\par\medskip

\subsection*{H-041 \quad Respuesta C}\label{s:H-041}
\textbf{Razonamiento.} Por Amdahl: 1/(0.4+0.6/4)=1/0.55≈1.82.

\textbf{A) Distractor. }El 40\% restante limita la mejora.

\textbf{B) Distractor. }Ignora el tiempo aún consumido por la parte acelerada.

\textbf{Repaso:} \hyperref[sec:h2]{Arquitectura de computadoras y su organización}. \textbf{Fuentes:} \cite{intel,architecture}. \hyperref[q:H-041]{Volver al reactivo}.\par\medskip

\subsection*{H-042 \quad Respuesta A}\label{s:H-042}
\textbf{Razonamiento.} El tiempo depende del producto de instrucciones y CPI dividido por frecuencia.

\textbf{B) Distractor. }El color no determina el trabajo arquitectónico.

\textbf{C) Distractor. }El gabinete por sí solo no determina el tiempo de ejecución.

\textbf{Repaso:} \hyperref[sec:h2]{Arquitectura de computadoras y su organización}. \textbf{Fuentes:} \cite{intel,architecture}. \hyperref[q:H-042]{Volver al reactivo}.\par\medskip

\subsection*{H-043 \quad Respuesta C}\label{s:H-043}
\textbf{Razonamiento.} Un plazo exige acotar la respuesta máxima, no solo su promedio.

\textbf{A) Distractor. }La capacidad no demuestra cumplimiento temporal.

\textbf{B) Distractor. }La pantalla no establece la latencia del servicio.

\textbf{Repaso:} \hyperref[sec:h2]{Arquitectura de computadoras y su organización}. \textbf{Fuentes:} \cite{intel,architecture}. \hyperref[q:H-043]{Volver al reactivo}.\par\medskip

\subsection*{H-044 \quad Respuesta A}\label{s:H-044}
\textbf{Razonamiento.} El recurso compartido puede crear un cuello de botella en el acceso.

\textbf{B) Distractor. }La representación de enteros no explica la contención.

\textbf{C) Distractor. }Esta organización no impide implementar interrupciones.

\textbf{Repaso:} \hyperref[sec:h2]{Arquitectura de computadoras y su organización}. \textbf{Fuentes:} \cite{intel,architecture}. \hyperref[q:H-044]{Volver al reactivo}.\par\medskip

\section{Sistemas embebidos}

\subsection*{H-045 \quad Respuesta A}\label{s:H-045}
\textbf{Razonamiento.} La integración permite controlar dispositivos con pocos componentes externos.

\textbf{B) Distractor. }No incluye la capacidad de ejecutar el control descrito.

\textbf{C) Distractor. }Carece del procesamiento programable requerido.

\textbf{Repaso:} \hyperref[sec:h3]{Sistemas embebidos}. \textbf{Fuentes:} \cite{avr,arduino,sensors}. \hyperref[q:H-045]{Volver al reactivo}.\par\medskip

\subsection*{H-046 \quad Respuesta B}\label{s:H-046}
\textbf{Razonamiento.} El firmware inicializa y gobierna el hardware y sus funciones.

\textbf{A) Distractor. }No es una categoría de protocolo de transmisión.

\textbf{C) Distractor. }La carcasa es un componente físico.

\textbf{Repaso:} \hyperref[sec:h3]{Sistemas embebidos}. \textbf{Fuentes:} \cite{avr,arduino,sensors}. \hyperref[q:H-046]{Volver al reactivo}.\par\medskip

\subsection*{H-047 \quad Respuesta A}\label{s:H-047}
\textbf{Razonamiento.} Un temporizador de vigilancia detecta falta de progreso mediante expiración.

\textbf{B) Distractor. }No modifica la capacidad de memoria.

\textbf{C) Distractor. }Mitiga bloqueos, pero no prueba corrección total.

\textbf{Repaso:} \hyperref[sec:h3]{Sistemas embebidos}. \textbf{Fuentes:} \cite{avr,arduino,sensors}. \hyperref[q:H-047]{Volver al reactivo}.\par\medskip

\subsection*{H-048 \quad Respuesta B}\label{s:H-048}
\textbf{Razonamiento.} I²C emplea datos y reloj compartidos con líneas de drenador/colector abierto.

\textbf{A) Distractor. }UART no comparte una línea de reloj de ese tipo.

\textbf{C) Distractor. }SPI usa señales diferenciadas y selección de esclavo.

\textbf{Repaso:} \hyperref[sec:h3]{Sistemas embebidos}. \textbf{Fuentes:} \cite{avr,arduino,sensors}. \hyperref[q:H-048]{Volver al reactivo}.\par\medskip

\subsection*{H-049 \quad Respuesta C}\label{s:H-049}
\textbf{Razonamiento.} El convertidor analógico-digital cuantifica una señal de entrada.

\textbf{A) Distractor. }El DAC hace la conversión inversa.

\textbf{B) Distractor. }PWM genera pulsos y no mide una tensión.

\textbf{Repaso:} \hyperref[sec:h3]{Sistemas embebidos}. \textbf{Fuentes:} \cite{avr,arduino,sensors}. \hyperref[q:H-049]{Volver al reactivo}.\par\medskip

\subsection*{H-050 \quad Respuesta B}\label{s:H-050}
\textbf{Razonamiento.} La CPU recibe una solicitud de atención en vez de depender solo de sondeo continuo.

\textbf{A) Distractor. }El mecanismo requiere energía como el resto del sistema.

\textbf{C) Distractor. }La rutina de servicio sigue ejecutando instrucciones.

\textbf{Repaso:} \hyperref[sec:h3]{Sistemas embebidos}. \textbf{Fuentes:} \cite{avr,arduino,sensors}. \hyperref[q:H-050]{Volver al reactivo}.\par\medskip

\subsection*{H-051 \quad Respuesta C}\label{s:H-051}
\textbf{Razonamiento.} El giroscopio entrega una medida relacionada con la rotación por unidad de tiempo.

\textbf{A) Distractor. }El magnetómetro mide campo magnético.

\textbf{B) Distractor. }El termistor responde a temperatura.

\textbf{Repaso:} \hyperref[sec:h3]{Sistemas embebidos}. \textbf{Fuentes:} \cite{avr,arduino,sensors}. \hyperref[q:H-051]{Volver al reactivo}.\par\medskip

\subsection*{H-052 \quad Respuesta A}\label{s:H-052}
\textbf{Razonamiento.} Importan el resultado y el instante en que se entrega.

\textbf{B) Distractor. }La frecuencia no demuestra cumplimiento de plazos.

\textbf{C) Distractor. }Los sistemas de tiempo real pueden ser programables.

\textbf{Repaso:} \hyperref[sec:h3]{Sistemas embebidos}. \textbf{Fuentes:} \cite{avr,arduino,sensors}. \hyperref[q:H-052]{Volver al reactivo}.\par\medskip

\subsection*{H-053 \quad Respuesta C}\label{s:H-053}
\textbf{Razonamiento.} 5/1024=0.0048828 V.

\textbf{A) Distractor. }Corresponde a un error decimal de un orden.

\textbf{B) Distractor. }No concuerda con los 1024 intervalos.

\textbf{Repaso:} \hyperref[sec:h3]{Sistemas embebidos}. \textbf{Fuentes:} \cite{avr,arduino,sensors}. \hyperref[q:H-053]{Volver al reactivo}.\par\medskip

\subsection*{H-054 \quad Respuesta A}\label{s:H-054}
\textbf{Razonamiento.} D=Ton/T=0.5/2=0.25.

\textbf{B) Distractor. }Requeriría 1 ms en alto.

\textbf{C) Distractor. }Es la fracción del periodo en bajo.

\textbf{Repaso:} \hyperref[sec:h3]{Sistemas embebidos}. \textbf{Fuentes:} \cite{avr,arduino,sensors}. \hyperref[q:H-054]{Volver al reactivo}.\par\medskip

\subsection*{H-055 \quad Respuesta B}\label{s:H-055}
\textbf{Razonamiento.} 1000/1000000=0.001 s.

\textbf{A) Distractor. }Es el periodo de un solo pulso.

\textbf{C) Distractor. }Requeriría un millón de pulsos.

\textbf{Repaso:} \hyperref[sec:h3]{Sistemas embebidos}. \textbf{Fuentes:} \cite{avr,arduino,sensors}. \hyperref[q:H-055]{Volver al reactivo}.\par\medskip

\subsection*{H-056 \quad Respuesta A}\label{s:H-056}
\textbf{Razonamiento.} Hay diez bits por trama: 10/9600 s.

\textbf{B) Distractor. }Cuenta datos pero omite inicio y parada.

\textbf{C) Distractor. }Invierte la relación entre tasa y tiempo.

\textbf{Repaso:} \hyperref[sec:h3]{Sistemas embebidos}. \textbf{Fuentes:} \cite{avr,arduino,sensors}. \hyperref[q:H-056]{Volver al reactivo}.\par\medskip

\subsection*{H-057 \quad Respuesta B}\label{s:H-057}
\textbf{Razonamiento.} La validación temporal evita contar rebotes como pulsaciones distintas.

\textbf{A) Distractor. }El endianness no elimina el ruido de contacto.

\textbf{C) Distractor. }Contaría varios eventos por una acción.

\textbf{Repaso:} \hyperref[sec:h3]{Sistemas embebidos}. \textbf{Fuentes:} \cite{avr,arduino,sensors}. \hyperref[q:H-057]{Volver al reactivo}.\par\medskip

\subsection*{H-058 \quad Respuesta C}\label{s:H-058}
\textbf{Razonamiento.} La ISR puede modificar el valor entre ambas lecturas; se necesita una captura atómica apropiada.

\textbf{A) Distractor. }El ancho de la CPU no cambia por el tipo de variable.

\textbf{B) Distractor. }Sí se pueden usar con sincronización.

\textbf{Repaso:} \hyperref[sec:h3]{Sistemas embebidos}. \textbf{Fuentes:} \cite{avr,arduino,sensors}. \hyperref[q:H-058]{Volver al reactivo}.\par\medskip

\subsection*{H-059 \quad Respuesta B}\label{s:H-059}
\textbf{Razonamiento.} volatile obliga a respetar accesos observables según las reglas del lenguaje y plataforma.

\textbf{A) Distractor. }No elimina carreras entre tareas o interrupciones.

\textbf{C) Distractor. }No une automáticamente múltiples instrucciones.

\textbf{Repaso:} \hyperref[sec:h3]{Sistemas embebidos}. \textbf{Fuentes:} \cite{avr,arduino,sensors}. \hyperref[q:H-059]{Volver al reactivo}.\par\medskip

\subsection*{H-060 \quad Respuesta C}\label{s:H-060}
\textbf{Razonamiento.} Una línea de selección activa identifica el esclavo correspondiente.

\textbf{A) Distractor. }SPI no requiere direccionamiento IP.

\textbf{B) Distractor. }El bit de parada pertenece a otra interfaz.

\textbf{Repaso:} \hyperref[sec:h3]{Sistemas embebidos}. \textbf{Fuentes:} \cite{avr,arduino,sensors}. \hyperref[q:H-060]{Volver al reactivo}.\par\medskip

\subsection*{H-061 \quad Respuesta A}\label{s:H-061}
\textbf{Razonamiento.} La lectura depende de orientación y aceleración específica; reposo no implica cero en todos los ejes.

\textbf{B) Distractor. }La medida es esperable bajo esas condiciones.

\textbf{C) Distractor. }La posición geográfica no es su magnitud de salida.

\textbf{Repaso:} \hyperref[sec:h3]{Sistemas embebidos}. \textbf{Fuentes:} \cite{avr,arduino,sensors}. \hyperref[q:H-061]{Volver al reactivo}.\par\medskip

\subsection*{H-062 \quad Respuesta C}\label{s:H-062}
\textbf{Razonamiento.} Se deben adaptar escala, impedancia y límites eléctricos.

\textbf{A) Distractor. }Puede sobrepasar los máximos permitidos.

\textbf{B) Distractor. }El tipo de dato no protege el pin físico.

\textbf{Repaso:} \hyperref[sec:h3]{Sistemas embebidos}. \textbf{Fuentes:} \cite{avr,arduino,sensors}. \hyperref[q:H-062]{Volver al reactivo}.\par\medskip

\subsection*{H-063 \quad Respuesta A}\label{s:H-063}
\textbf{Razonamiento.} Primero se selecciona la entrada; después se convierte, se espera y se recupera el dato.

\textbf{B) Distractor. }Lee antes de efectuar la conversión requerida.

\textbf{C) Distractor. }Lee antes de completar la adquisición y configura demasiado tarde.

\textbf{Repaso:} \hyperref[sec:h3]{Sistemas embebidos}. \textbf{Fuentes:} \cite{avr,arduino,sensors}. \hyperref[q:H-063]{Volver al reactivo}.\par\medskip

\subsection*{H-064 \quad Respuesta B}\label{s:H-064}
\textbf{Razonamiento.} Cada dispositivo observa una magnitud diferente que puede fusionarse con otras.

\textbf{A) Distractor. }No corresponde a los principios de medida.

\textbf{C) Distractor. }Intercambia posición, campo y rotación.

\textbf{Repaso:} \hyperref[sec:h3]{Sistemas embebidos}. \textbf{Fuentes:} \cite{avr,arduino,sensors}. \hyperref[q:H-064]{Volver al reactivo}.\par\medskip

\subsection*{H-065 \quad Respuesta A}\label{s:H-065}
\textbf{Razonamiento.} Se exige fs>2×3 kHz; 8 kHz cumple.

\textbf{B) Distractor. }La igualdad no supera estrictamente el límite y es un caso fronterizo.

\textbf{C) Distractor. }Está por debajo del doble del ancho de banda.

\textbf{Repaso:} \hyperref[sec:h3]{Sistemas embebidos}. \textbf{Fuentes:} \cite{avr,arduino,sensors}. \hyperref[q:H-065]{Volver al reactivo}.\par\medskip

\subsection*{H-066 \quad Respuesta B}\label{s:H-066}
\textbf{Razonamiento.} Se pueden acumular o perder eventos; debe reducirse el trabajo o cambiarse el diseño.

\textbf{A) Distractor. }Un exceso de servicio no garantiza temporización estable.

\textbf{C) Distractor. }El tiempo de servicio no añade bits al ADC.

\textbf{Repaso:} \hyperref[sec:h3]{Sistemas embebidos}. \textbf{Fuentes:} \cite{avr,arduino,sensors}. \hyperref[q:H-066]{Volver al reactivo}.\par\medskip

\subsection*{H-067 \quad Respuesta C}\label{s:H-067}
\textbf{Razonamiento.} Un transistor o controlador permite entregar corriente sin sobrecargar el GPIO.

\textbf{A) Distractor. }Puede dañar la salida.

\textbf{B) Distractor. }Las entradas no proporcionan potencia al motor.

\textbf{Repaso:} \hyperref[sec:h3]{Sistemas embebidos}. \textbf{Fuentes:} \cite{avr,arduino,sensors}. \hyperref[q:H-067]{Volver al reactivo}.\par\medskip

\subsection*{H-068 \quad Respuesta B}\label{s:H-068}
\textbf{Razonamiento.} El reposo reduce actividad y un evento habilitado reactiva el procesamiento.

\textbf{A) Distractor. }Puede aumentar consumo sin resolver la espera.

\textbf{C) Distractor. }Sigue consumiendo ciclos durante la espera.

\textbf{Repaso:} \hyperref[sec:h3]{Sistemas embebidos}. \textbf{Fuentes:} \cite{avr,arduino,sensors}. \hyperref[q:H-068]{Volver al reactivo}.\par\medskip

\subsection*{H-069 \quad Respuesta C}\label{s:H-069}
\textbf{Razonamiento.} La herencia limita la inversión de prioridad mientras se libera el recurso.

\textbf{A) Distractor. }Empeora la capacidad de cumplir plazos.

\textbf{B) Distractor. }Cambiar el requisito no corrige el comportamiento requerido.

\textbf{Repaso:} \hyperref[sec:h3]{Sistemas embebidos}. \textbf{Fuentes:} \cite{avr,arduino,sensors}. \hyperref[q:H-069]{Volver al reactivo}.\par\medskip

\subsection*{H-070 \quad Respuesta A}\label{s:H-070}
\textbf{Razonamiento.} La validación de límites impide escribir más allá del espacio asignado.

\textbf{B) Distractor. }El receptor debe tolerar entradas malformadas.

\textbf{C) Distractor. }No evita una copia fuera de límites.

\textbf{Repaso:} \hyperref[sec:h3]{Sistemas embebidos}. \textbf{Fuentes:} \cite{avr,arduino,sensors}. \hyperref[q:H-070]{Volver al reactivo}.\par\medskip

\subsection*{H-071 \quad Respuesta C}\label{s:H-071}
\textbf{Razonamiento.} La validación y una ruta de recuperación evitan depender de una imagen incompleta.

\textbf{A) Distractor. }Puede destruir precisamente el mecanismo de recuperación.

\textbf{B) Distractor. }Aumenta el riesgo de ejecutar contenido corrupto.

\textbf{Repaso:} \hyperref[sec:h3]{Sistemas embebidos}. \textbf{Fuentes:} \cite{avr,arduino,sensors}. \hyperref[q:H-071]{Volver al reactivo}.\par\medskip

\section{Automatización y control de procesos}

\subsection*{H-072 \quad Respuesta A}\label{s:H-072}
\textbf{Razonamiento.} La salida medida influye en la acción de control.

\textbf{B) Distractor. }La realimentación necesita información de la variable.

\textbf{C) Distractor. }La estabilidad depende de dinámica y parámetros.

\textbf{Repaso:} \hyperref[sec:h4]{Automatización y control de procesos}. \textbf{Fuentes:} \cite{control,signals,pid}. \hyperref[q:H-072]{Volver al reactivo}.\par\medskip

\subsection*{H-073 \quad Respuesta B}\label{s:H-073}
\textbf{Razonamiento.} e=r−y compara el objetivo con el resultado.

\textbf{A) Distractor. }El tiempo no es la diferencia de seguimiento.

\textbf{C) Distractor. }Corresponde a otro signo de realimentación.

\textbf{Repaso:} \hyperref[sec:h4]{Automatización y control de procesos}. \textbf{Fuentes:} \cite{control,signals,pid}. \hyperref[q:H-073]{Volver al reactivo}.\par\medskip

\subsection*{H-074 \quad Respuesta A}\label{s:H-074}
\textbf{Razonamiento.} G(s)=Y(s)/U(s) representa la dinámica entrada-salida bajo esas condiciones.

\textbf{B) Distractor. }No se limita a resistencias eléctricas.

\textbf{C) Distractor. }No es una operación estadística genérica.

\textbf{Repaso:} \hyperref[sec:h4]{Automatización y control de procesos}. \textbf{Fuentes:} \cite{control,signals,pid}. \hyperref[q:H-074]{Volver al reactivo}.\par\medskip

\subsection*{H-075 \quad Respuesta B}\label{s:H-075}
\textbf{Razonamiento.} La integral considera el área temporal del error.

\textbf{A) Distractor. }La derivativa usa la tasa de cambio.

\textbf{C) Distractor. }La proporcional usa el error instantáneo.

\textbf{Repaso:} \hyperref[sec:h4]{Automatización y control de procesos}. \textbf{Fuentes:} \cite{control,signals,pid}. \hyperref[q:H-075]{Volver al reactivo}.\par\medskip

\subsection*{H-076 \quad Respuesta C}\label{s:H-076}
\textbf{Razonamiento.} El actuador transforma energía para modificar el proceso.

\textbf{A) Distractor. }El sensor mide la variable.

\textbf{B) Distractor. }Un archivo almacena información pero no actúa físicamente.

\textbf{Repaso:} \hyperref[sec:h4]{Automatización y control de procesos}. \textbf{Fuentes:} \cite{control,signals,pid}. \hyperref[q:H-076]{Volver al reactivo}.\par\medskip

\subsection*{H-077 \quad Respuesta B}\label{s:H-077}
\textbf{Razonamiento.} Se cumplen aditividad y homogeneidad.

\textbf{A) Distractor. }No exige periodicidad de las señales.

\textbf{C) Distractor. }La linealidad no exige salida constante.

\textbf{Repaso:} \hyperref[sec:h4]{Automatización y control de procesos}. \textbf{Fuentes:} \cite{control,signals,pid}. \hyperref[q:H-077]{Volver al reactivo}.\par\medskip

\subsection*{H-078 \quad Respuesta C}\label{s:H-078}
\textbf{Razonamiento.} La dinámica no cambia al trasladar el experimento en el tiempo.

\textbf{A) Distractor. }Puede generar salidas no nulas.

\textbf{B) Distractor. }Puede depender de valores pasados.

\textbf{Repaso:} \hyperref[sec:h4]{Automatización y control de procesos}. \textbf{Fuentes:} \cite{control,signals,pid}. \hyperref[q:H-078]{Volver al reactivo}.\par\medskip

\subsection*{H-079 \quad Respuesta A}\label{s:H-079}
\textbf{Razonamiento.} Se requiere transducir temperatura a información utilizable.

\textbf{B) Distractor. }El contactor conmuta potencia, pero por sí solo no mide.

\textbf{C) Distractor. }El engrane transmite movimiento.

\textbf{Repaso:} \hyperref[sec:h4]{Automatización y control de procesos}. \textbf{Fuentes:} \cite{control,signals,pid}. \hyperref[q:H-079]{Volver al reactivo}.\par\medskip

\subsection*{H-080 \quad Respuesta C}\label{s:H-080}
\textbf{Razonamiento.} u=Kp×e=3×2=6.

\textbf{A) Distractor. }Divide en vez de multiplicar.

\textbf{B) Distractor. }Suma magnitudes en lugar de aplicar la ley proporcional.

\textbf{Repaso:} \hyperref[sec:h4]{Automatización y control de procesos}. \textbf{Fuentes:} \cite{control,signals,pid}. \hyperref[q:H-080]{Volver al reactivo}.\par\medskip

\subsection*{H-081 \quad Respuesta A}\label{s:H-081}
\textbf{Razonamiento.} La fracción alcanzada a τ es 1−e\textasciicircum{}(−1)≈0.632.

\textbf{B) Distractor. }La aproximación al valor final es asintótica.

\textbf{C) Distractor. }Es la fracción aún faltante.

\textbf{Repaso:} \hyperref[sec:h4]{Automatización y control de procesos}. \textbf{Fuentes:} \cite{control,signals,pid}. \hyperref[q:H-081]{Volver al reactivo}.\par\medskip

\subsection*{H-082 \quad Respuesta B}\label{s:H-082}
\textbf{Razonamiento.} La fracción de escala es (12−4)/(20−4)=0.5.

\textbf{A) Distractor. }Confunde la unidad eléctrica con la temperatura.

\textbf{C) Distractor. }No descuenta correctamente el cero vivo de 4 mA.

\textbf{Repaso:} \hyperref[sec:h4]{Automatización y control de procesos}. \textbf{Fuentes:} \cite{control,signals,pid}. \hyperref[q:H-082]{Volver al reactivo}.\par\medskip

\subsection*{H-083 \quad Respuesta A}\label{s:H-083}
\textbf{Razonamiento.} Los dos umbrales evitan cambios repetidos cerca de un solo punto.

\textbf{B) Distractor. }No cambia las propiedades físicas del proceso.

\textbf{C) Distractor. }La salida sigue siendo binaria.

\textbf{Repaso:} \hyperref[sec:h4]{Automatización y control de procesos}. \textbf{Fuentes:} \cite{control,signals,pid}. \hyperref[q:H-083]{Volver al reactivo}.\par\medskip

\subsection*{H-084 \quad Respuesta B}\label{s:H-084}
\textbf{Razonamiento.} El polo es la raíz del denominador s+4=0.

\textbf{A) Distractor. }El denominador no se anula en cero.

\textbf{C) Distractor. }Invierte el signo al resolver.

\textbf{Repaso:} \hyperref[sec:h4]{Automatización y control de procesos}. \textbf{Fuentes:} \cite{control,signals,pid}. \hyperref[q:H-084]{Volver al reactivo}.\par\medskip

\subsection*{H-085 \quad Respuesta C}\label{s:H-085}
\textbf{Razonamiento.} Los modos asociados decaen exponencialmente.

\textbf{A) Distractor. }Los polos con parte real negativa no producen crecimiento libre.

\textbf{B) Distractor. }La estabilidad no elimina la influencia de la entrada.

\textbf{Repaso:} \hyperref[sec:h4]{Automatización y control de procesos}. \textbf{Fuentes:} \cite{control,signals,pid}. \hyperref[q:H-085]{Volver al reactivo}.\par\medskip

\subsection*{H-086 \quad Respuesta B}\label{s:H-086}
\textbf{Razonamiento.} Un pasa-bajas reduce componentes rápidas y conserva la banda lenta de interés.

\textbf{A) Distractor. }No mejora la relación entre señal lenta y ruido rápido.

\textbf{C) Distractor. }Favorece precisamente la región que se desea atenuar.

\textbf{Repaso:} \hyperref[sec:h4]{Automatización y control de procesos}. \textbf{Fuentes:} \cite{control,signals,pid}. \hyperref[q:H-086]{Volver al reactivo}.\par\medskip

\subsection*{H-087 \quad Respuesta C}\label{s:H-087}
\textbf{Razonamiento.} La multiplicación por una constante conserva superposición y usa solo el instante actual.

\textbf{A) Distractor. }Una ganancia constante es una operación lineal.

\textbf{B) Distractor. }No aparece ningún valor futuro.

\textbf{Repaso:} \hyperref[sec:h4]{Automatización y control de procesos}. \textbf{Fuentes:} \cite{control,signals,pid}. \hyperref[q:H-087]{Volver al reactivo}.\par\medskip

\subsection*{H-088 \quad Respuesta A}\label{s:H-088}
\textbf{Razonamiento.} Viola homogeneidad: T(2u)=4u², distinto de 2T(u).

\textbf{B) Distractor. }Un sistema lineal también puede tener salidas positivas.

\textbf{C) Distractor. }La continuidad no determina linealidad.

\textbf{Repaso:} \hyperref[sec:h4]{Automatización y control de procesos}. \textbf{Fuentes:} \cite{control,signals,pid}. \hyperref[q:H-088]{Volver al reactivo}.\par\medskip

\subsection*{H-089 \quad Respuesta C}\label{s:H-089}
\textbf{Razonamiento.} Ts=1/fs=0.01 s.

\textbf{A) Distractor. }Equivale a una frecuencia de 10 Hz.

\textbf{B) Distractor. }Equivale a 1000 Hz.

\textbf{Repaso:} \hyperref[sec:h4]{Automatización y control de procesos}. \textbf{Fuentes:} \cite{control,signals,pid}. \hyperref[q:H-089]{Volver al reactivo}.\par\medskip

\subsection*{H-090 \quad Respuesta A}\label{s:H-090}
\textbf{Razonamiento.} Se suman las tres muestras y se divide entre tres: 18/3=6.

\textbf{B) Distractor. }Es la suma, no la media.

\textbf{C) Distractor. }Es solo una de las muestras.

\textbf{Repaso:} \hyperref[sec:h4]{Automatización y control de procesos}. \textbf{Fuentes:} \cite{control,signals,pid}. \hyperref[q:H-090]{Volver al reactivo}.\par\medskip

\subsection*{H-091 \quad Respuesta B}\label{s:H-091}
\textbf{Razonamiento.} La acumulación durante saturación puede retrasar la recuperación y causar sobreimpulso.

\textbf{A) Distractor. }No describe la saturación del control.

\textbf{C) Distractor. }Nyquist trata de muestreo o estabilidad según contexto, no de esta acumulación.

\textbf{Repaso:} \hyperref[sec:h4]{Automatización y control de procesos}. \textbf{Fuentes:} \cite{control,signals,pid}. \hyperref[q:H-091]{Volver al reactivo}.\par\medskip

\subsection*{H-092 \quad Respuesta A}\label{s:H-092}
\textbf{Razonamiento.} Se adquiere el estado del proceso, se calcula la orden y se actualiza la salida.

\textbf{B) Distractor. }Actúa con información aún no calculada.

\textbf{C) Distractor. }Calcula el error antes de disponer de la medición actual.

\textbf{Repaso:} \hyperref[sec:h4]{Automatización y control de procesos}. \textbf{Fuentes:} \cite{control,signals,pid}. \hyperref[q:H-092]{Volver al reactivo}.\par\medskip

\subsection*{H-093 \quad Respuesta B}\label{s:H-093}
\textbf{Razonamiento.} Las tres acciones usan aspectos diferentes del error.

\textbf{A) Distractor. }No corresponde a las expresiones de un PID.

\textbf{C) Distractor. }Intercambia las definiciones de las acciones.

\textbf{Repaso:} \hyperref[sec:h4]{Automatización y control de procesos}. \textbf{Fuentes:} \cite{control,signals,pid}. \hyperref[q:H-093]{Volver al reactivo}.\par\medskip

\subsection*{H-094 \quad Respuesta C}\label{s:H-094}
\textbf{Razonamiento.} T=G/(1+G)=2/(s+1+2)=2/(s+3).

\textbf{A) Distractor. }Usa incorrectamente el signo de realimentación.

\textbf{B) Distractor. }Es la planta sin cerrar el lazo.

\textbf{Repaso:} \hyperref[sec:h4]{Automatización y control de procesos}. \textbf{Fuentes:} \cite{control,signals,pid}. \hyperref[q:H-094]{Volver al reactivo}.\par\medskip

\subsection*{H-095 \quad Respuesta B}\label{s:H-095}
\textbf{Razonamiento.} La realimentación puede volverse poco amortiguada o inestable con ganancia excesiva.

\textbf{A) Distractor. }Sin medición no se corrige el error de salida del mismo modo.

\textbf{C) Distractor. }La mejora no es monótona para cualquier planta.

\textbf{Repaso:} \hyperref[sec:h4]{Automatización y control de procesos}. \textbf{Fuentes:} \cite{control,signals,pid}. \hyperref[q:H-095]{Volver al reactivo}.\par\medskip

\subsection*{H-096 \quad Respuesta C}\label{s:H-096}
\textbf{Razonamiento.} La demora añade fase y variación temporal que pueden deteriorar el lazo.

\textbf{A) Distractor. }El estilo visual no explica la dinámica de control.

\textbf{B) Distractor. }No influye en la respuesta temporal del proceso.

\textbf{Repaso:} \hyperref[sec:h4]{Automatización y control de procesos}. \textbf{Fuentes:} \cite{control,signals,pid}. \hyperref[q:H-096]{Volver al reactivo}.\par\medskip

\subsection*{H-097 \quad Respuesta A}\label{s:H-097}
\textbf{Razonamiento.} El sesgo es un error sistemático corregible mediante calibración.

\textbf{B) Distractor. }Mostrar más cifras no elimina el desplazamiento.

\textbf{C) Distractor. }La derivada de un sesgo constante no corrige la medida absoluta.

\textbf{Repaso:} \hyperref[sec:h4]{Automatización y control de procesos}. \textbf{Fuentes:} \cite{control,signals,pid}. \hyperref[q:H-097]{Volver al reactivo}.\par\medskip

\subsection*{H-098 \quad Respuesta C}\label{s:H-098}
\textbf{Razonamiento.} La ley de control debe respetar los estados y tiempos del actuador.

\textbf{A) Distractor. }La orden no será realizable por ese hardware.

\textbf{B) Distractor. }Contradice la válvula binaria especificada.

\textbf{Repaso:} \hyperref[sec:h4]{Automatización y control de procesos}. \textbf{Fuentes:} \cite{control,signals,pid}. \hyperref[q:H-098]{Volver al reactivo}.\par\medskip

\subsection*{H-099 \quad Respuesta A}\label{s:H-099}
\textbf{Razonamiento.} La medición más suave puede introducir fase y perjudicar la respuesta.

\textbf{B) Distractor. }No considera los efectos dinámicos del filtro.

\textbf{C) Distractor. }Los lazos son sensibles al retardo.

\textbf{Repaso:} \hyperref[sec:h4]{Automatización y control de procesos}. \textbf{Fuentes:} \cite{control,signals,pid}. \hyperref[q:H-099]{Volver al reactivo}.\par\medskip

\subsection*{H-100 \quad Respuesta B}\label{s:H-100}
\textbf{Razonamiento.} La rotación, aceleración específica y campo aportan información distinta para estimar orientación.

\textbf{A) Distractor. }La fusión no garantiza eliminar sesgos sin calibración.

\textbf{C) Distractor. }Las magnitudes físicas y fuentes de error son diferentes.

\textbf{Repaso:} \hyperref[sec:h4]{Automatización y control de procesos}. \textbf{Fuentes:} \cite{control,signals,pid}. \hyperref[q:H-100]{Volver al reactivo}.\par\medskip

\section{Diseño de redes}

\subsection*{R-001 \quad Respuesta C}\label{s:R-001}
\textbf{Razonamiento.} La capa de red transporta paquetes entre redes mediante direccionamiento lógico.

\textbf{A) Distractor. }La sesión organiza el diálogo entre aplicaciones.

\textbf{B) Distractor. }La presentación se relaciona con representación de datos.

\textbf{Repaso:} \hyperref[sec:r1]{Diseño de redes}. \textbf{Fuentes:} \cite{rfc8200,rfc4632,networks,vlan}. \hyperref[q:R-001]{Volver al reactivo}.\par\medskip

\subsection*{R-002 \quad Respuesta A}\label{s:R-002}
\textbf{Razonamiento.} El switch aprende MAC de origen y decide puertos de salida.

\textbf{B) Distractor. }El repetidor regenera señales sin esa tabla.

\textbf{C) Distractor. }No es un dispositivo de conmutación Ethernet.

\textbf{Repaso:} \hyperref[sec:r1]{Diseño de redes}. \textbf{Fuentes:} \cite{rfc8200,rfc4632,networks,vlan}. \hyperref[q:R-002]{Volver al reactivo}.\par\medskip

\subsection*{R-003 \quad Respuesta C}\label{s:R-003}
\textbf{Razonamiento.} IPv6 emplea direcciones de 128 bits.

\textbf{A) Distractor. }Es la longitud de una dirección IPv4.

\textbf{B) Distractor. }Es la longitud usual de una MAC Ethernet.

\textbf{Repaso:} \hyperref[sec:r1]{Diseño de redes}. \textbf{Fuentes:} \cite{rfc8200,rfc4632,networks,vlan}. \hyperref[q:R-003]{Volver al reactivo}.\par\medskip

\subsection*{R-004 \quad Respuesta A}\label{s:R-004}
\textbf{Razonamiento.} La fibra transporta luz y no capta interferencia electromagnética como un conductor.

\textbf{B) Distractor. }El cobre puede sufrir interferencia.

\textbf{C) Distractor. }El coaxial es conductor y no carece de pérdidas.

\textbf{Repaso:} \hyperref[sec:r1]{Diseño de redes}. \textbf{Fuentes:} \cite{rfc8200,rfc4632,networks,vlan}. \hyperref[q:R-004]{Volver al reactivo}.\par\medskip

\subsection*{R-005 \quad Respuesta B}\label{s:R-005}
\textbf{Razonamiento.} DNS consulta registros como A y AAAA.

\textbf{A) Distractor. }NTP sincroniza tiempo.

\textbf{C) Distractor. }DHCP asigna configuración; no sustituye el sistema de nombres.

\textbf{Repaso:} \hyperref[sec:r1]{Diseño de redes}. \textbf{Fuentes:} \cite{rfc8200,rfc4632,networks,vlan}. \hyperref[q:R-005]{Volver al reactivo}.\par\medskip

\subsection*{R-006 \quad Respuesta A}\label{s:R-006}
\textbf{Razonamiento.} La longitud del prefijo fija cuántos de los 32 bits pertenecen a la red.

\textbf{B) Distractor. }Toda la dirección IPv4 tiene cuatro bytes.

\textbf{C) Distractor. }El número de hosts depende de los bits restantes.

\textbf{Repaso:} \hyperref[sec:r1]{Diseño de redes}. \textbf{Fuentes:} \cite{rfc8200,rfc4632,networks,vlan}. \hyperref[q:R-006]{Volver al reactivo}.\par\medskip

\subsection*{R-007 \quad Respuesta B}\label{s:R-007}
\textbf{Razonamiento.} La segmentación separa tráfico de difusión entre grupos lógicos.

\textbf{A) Distractor. }No obliga a usar direcciones públicas.

\textbf{C) Distractor. }La separación no cifra por sí misma.

\textbf{Repaso:} \hyperref[sec:r1]{Diseño de redes}. \textbf{Fuentes:} \cite{rfc8200,rfc4632,networks,vlan}. \hyperref[q:R-007]{Volver al reactivo}.\par\medskip

\subsection*{R-008 \quad Respuesta C}\label{s:R-008}
\textbf{Razonamiento.} UDP proporciona datagramas sin las garantías de secuencia y retransmisión de TCP.

\textbf{A) Distractor. }El establecimiento de conexión descrito corresponde a TCP.

\textbf{B) Distractor. }La aplicación debe resolver pérdidas si lo necesita.

\textbf{Repaso:} \hyperref[sec:r1]{Diseño de redes}. \textbf{Fuentes:} \cite{rfc8200,rfc4632,networks,vlan}. \hyperref[q:R-008]{Volver al reactivo}.\par\medskip

\subsection*{R-009 \quad Respuesta B}\label{s:R-009}
\textbf{Razonamiento.} Con seis bits de host hay 64 direcciones y 62 utilizables.

\textbf{A) Distractor. }Ofrece 14 hosts convencionales.

\textbf{C) Distractor. }Ofrece 30 hosts convencionales.

\textbf{Repaso:} \hyperref[sec:r1]{Diseño de redes}. \textbf{Fuentes:} \cite{rfc8200,rfc4632,networks,vlan}. \hyperref[q:R-009]{Volver al reactivo}.\par\medskip

\subsection*{R-010 \quad Respuesta C}\label{s:R-010}
\textbf{Razonamiento.} Los bloques avanzan de 64 en 64; 77 pertenece al bloque 64–127.

\textbf{A) Distractor. }Corresponde al primer bloque /26.

\textbf{B) Distractor. }No es un límite de bloque /26.

\textbf{Repaso:} \hyperref[sec:r1]{Diseño de redes}. \textbf{Fuentes:} \cite{rfc8200,rfc4632,networks,vlan}. \hyperref[q:R-010]{Volver al reactivo}.\par\medskip

\subsection*{R-011 \quad Respuesta A}\label{s:R-011}
\textbf{Razonamiento.} El último valor del bloque de 64 direcciones es 127.

\textbf{B) Distractor. }Es el último host convencional.

\textbf{C) Distractor. }Es la dirección de red.

\textbf{Repaso:} \hyperref[sec:r1]{Diseño de redes}. \textbf{Fuentes:} \cite{rfc8200,rfc4632,networks,vlan}. \hyperref[q:R-011]{Volver al reactivo}.\par\medskip

\subsection*{R-012 \quad Respuesta C}\label{s:R-012}
\textbf{Razonamiento.} Quedan tres bits: 2³−2=6.

\textbf{A) Distractor. }Incluye red y broadcast.

\textbf{B) Distractor. }Resta más direcciones de las reservadas.

\textbf{Repaso:} \hyperref[sec:r1]{Diseño de redes}. \textbf{Fuentes:} \cite{rfc8200,rfc4632,networks,vlan}. \hyperref[q:R-012]{Volver al reactivo}.\par\medskip

\subsection*{R-013 \quad Respuesta A}\label{s:R-013}
\textbf{Razonamiento.} Comparten 23 bits y forman dos bloques /24 contiguos y alineados.

\textbf{B) Distractor. }Es más pequeño que una sola de las redes.

\textbf{C) Distractor. }Incluye también 10.0.2.0/24 y 10.0.3.0/24.

\textbf{Repaso:} \hyperref[sec:r1]{Diseño de redes}. \textbf{Fuentes:} \cite{rfc8200,rfc4632,networks,vlan}. \hyperref[q:R-013]{Volver al reactivo}.\par\medskip

\subsection*{R-014 \quad Respuesta B}\label{s:R-014}
\textbf{Razonamiento.} ARP resuelve la dirección del vecino de enlace al que se entregará la trama.

\textbf{A) Distractor. }El host no necesita conocer todas las MAC intermedias.

\textbf{C) Distractor. }La MAC remota no se conserva a través de routers.

\textbf{Repaso:} \hyperref[sec:r1]{Diseño de redes}. \textbf{Fuentes:} \cite{rfc8200,rfc4632,networks,vlan}. \hyperref[q:R-014]{Volver al reactivo}.\par\medskip

\subsection*{R-015 \quad Respuesta A}\label{s:R-015}
\textbf{Razonamiento.} Tiempo=bits/tasa=10\textasciicircum{}7/10\textasciicircum{}8=0.1 s.

\textbf{B) Distractor. }Invierte la relación y las escalas.

\textbf{C) Distractor. }No corresponde a la tasa indicada.

\textbf{Repaso:} \hyperref[sec:r1]{Diseño de redes}. \textbf{Fuentes:} \cite{rfc8200,rfc4632,networks,vlan}. \hyperref[q:R-015]{Volver al reactivo}.\par\medskip

\subsection*{R-016 \quad Respuesta B}\label{s:R-016}
\textbf{Razonamiento.} fe80::/10 identifica direcciones de enlace local.

\textbf{A) Distractor. }::1 es loopback.

\textbf{C) Distractor. }2001:db8::/32 está reservado para documentación.

\textbf{Repaso:} \hyperref[sec:r1]{Diseño de redes}. \textbf{Fuentes:} \cite{rfc8200,rfc4632,networks,vlan}. \hyperref[q:R-016]{Volver al reactivo}.\par\medskip

\subsection*{R-017 \quad Respuesta C}\label{s:R-017}
\textbf{Razonamiento.} El tráfico entre dominios IP requiere un router o switch de capa 3.

\textbf{A) Distractor. }Puede crear un bucle y no define el encaminamiento requerido.

\textbf{B) Distractor. }El tamaño de trama no habilita comunicación entre redes.

\textbf{Repaso:} \hyperref[sec:r1]{Diseño de redes}. \textbf{Fuentes:} \cite{rfc8200,rfc4632,networks,vlan}. \hyperref[q:R-017]{Volver al reactivo}.\par\medskip

\subsection*{R-018 \quad Respuesta B}\label{s:R-018}
\textbf{Razonamiento.} La etiqueta permite distinguir VLAN en el enlace compartido.

\textbf{A) Distractor. }No participa en el encabezado de red.

\textbf{C) Distractor. }DNS no autentica ni etiqueta tramas Ethernet.

\textbf{Repaso:} \hyperref[sec:r1]{Diseño de redes}. \textbf{Fuentes:} \cite{rfc8200,rfc4632,networks,vlan}. \hyperref[q:R-018]{Volver al reactivo}.\par\medskip

\subsection*{R-019 \quad Respuesta C}\label{s:R-019}
\textbf{Razonamiento.} Las necesidades guían la topología lógica, luego la selección física y la validación.

\textbf{A) Distractor. }Compra antes de conocer restricciones y valida sin diseño completo.

\textbf{B) Distractor. }La validación final necesita una propuesta basada en requisitos.

\textbf{Repaso:} \hyperref[sec:r1]{Diseño de redes}. \textbf{Fuentes:} \cite{rfc8200,rfc4632,networks,vlan}. \hyperref[q:R-019]{Volver al reactivo}.\par\medskip

\subsection*{R-020 \quad Respuesta A}\label{s:R-020}
\textbf{Razonamiento.} Cada capa opera sobre una abstracción distinta: señal, trama y paquete.

\textbf{B) Distractor. }Desplaza las funciones entre niveles.

\textbf{C) Distractor. }Confunde el direccionamiento lógico con el medio físico.

\textbf{Repaso:} \hyperref[sec:r1]{Diseño de redes}. \textbf{Fuentes:} \cite{rfc8200,rfc4632,networks,vlan}. \hyperref[q:R-020]{Volver al reactivo}.\par\medskip

\subsection*{R-021 \quad Respuesta C}\label{s:R-021}
\textbf{Razonamiento.} Se usa la coincidencia de prefijo más largo: /16.

\textbf{A) Distractor. }El criterio favorece la ruta más específica.

\textbf{B) Distractor. }Varias coincidencias son normales en enrutamiento.

\textbf{Repaso:} \hyperref[sec:r1]{Diseño de redes}. \textbf{Fuentes:} \cite{rfc8200,rfc4632,networks,vlan}. \hyperref[q:R-021]{Volver al reactivo}.\par\medskip

\subsection*{R-022 \quad Respuesta A}\label{s:R-022}
\textbf{Razonamiento.} STP construye una topología de reenvío sin ciclos.

\textbf{B) Distractor. }No es un protocolo de cifrado de aplicación.

\textbf{C) Distractor. }No administra el espacio de direcciones IPv6.

\textbf{Repaso:} \hyperref[sec:r1]{Diseño de redes}. \textbf{Fuentes:} \cite{rfc8200,rfc4632,networks,vlan}. \hyperref[q:R-022]{Volver al reactivo}.\par\medskip

\subsection*{R-023 \quad Respuesta B}\label{s:R-023}
\textbf{Razonamiento.} El encaminamiento convencional requiere distinguir redes de destino sin ambigüedad.

\textbf{A) Distractor. }La máscara es necesaria para decidir alcance local.

\textbf{C) Distractor. }Duplicar MAC introduce conflictos.

\textbf{Repaso:} \hyperref[sec:r1]{Diseño de redes}. \textbf{Fuentes:} \cite{rfc8200,rfc4632,networks,vlan}. \hyperref[q:R-023]{Volver al reactivo}.\par\medskip

\subsection*{R-024 \quad Respuesta A}\label{s:R-024}
\textbf{Razonamiento.} Los tiempos de llegada condicionan la interacción aunque la tasa media sea alta.

\textbf{B) Distractor. }No determina por sí sola el comportamiento temporal.

\textbf{C) Distractor. }La longitud del nombre no explica la variación de demora.

\textbf{Repaso:} \hyperref[sec:r1]{Diseño de redes}. \textbf{Fuentes:} \cite{rfc8200,rfc4632,networks,vlan}. \hyperref[q:R-024]{Volver al reactivo}.\par\medskip

\subsection*{R-025 \quad Respuesta B}\label{s:R-025}
\textbf{Razonamiento.} /25 da 126 hosts, /26 da 62 y /27 da 30; caben sin superposición.

\textbf{A) Distractor. }Un /28 solo aporta 14 hosts convencionales.

\textbf{C) Distractor. }Tres bloques de 128 direcciones exceden las 256 disponibles.

\textbf{Repaso:} \hyperref[sec:r1]{Diseño de redes}. \textbf{Fuentes:} \cite{rfc8200,rfc4632,networks,vlan}. \hyperref[q:R-025]{Volver al reactivo}.\par\medskip

\subsection*{R-026 \quad Respuesta C}\label{s:R-026}
\textbf{Razonamiento.} Una sola rotura puede inutilizar ambos enlaces aunque sean dos cables.

\textbf{A) Distractor. }El color no separa el riesgo físico.

\textbf{B) Distractor. }La etiqueta no elimina una causa común de falla.

\textbf{Repaso:} \hyperref[sec:r1]{Diseño de redes}. \textbf{Fuentes:} \cite{rfc8200,rfc4632,networks,vlan}. \hyperref[q:R-026]{Volver al reactivo}.\par\medskip

\subsection*{R-027 \quad Respuesta B}\label{s:R-027}
\textbf{Razonamiento.} La compatibilidad óptica y la pérdida total deben satisfacer ambos extremos.

\textbf{A) Distractor. }Hay límites por dispersión, atenuación y transceptores.

\textbf{C) Distractor. }La forma mecánica no garantiza compatibilidad ni potencia suficiente.

\textbf{Repaso:} \hyperref[sec:r1]{Diseño de redes}. \textbf{Fuentes:} \cite{rfc8200,rfc4632,networks,vlan}. \hyperref[q:R-027]{Volver al reactivo}.\par\medskip

\section{Administración de redes}

\subsection*{R-028 \quad Respuesta C}\label{s:R-028}
\textbf{Razonamiento.} DHCP entrega una concesión con dirección y otros parámetros.

\textbf{A) Distractor. }SMTP transporta correo.

\textbf{B) Distractor. }SSH proporciona acceso remoto seguro.

\textbf{Repaso:} \hyperref[sec:r2]{Administración de redes}. \textbf{Fuentes:} \cite{rfc2131,rfc1034,rfc9293,networks}. \hyperref[q:R-028]{Volver al reactivo}.\par\medskip

\subsection*{R-029 \quad Respuesta A}\label{s:R-029}
\textbf{Razonamiento.} La respuesta ICMP informa alcance y RTT bajo condiciones de la prueba.

\textbf{B) Distractor. }Un ping exitoso no demuestra salud de cada servicio.

\textbf{C) Distractor. }La medida eléctrica requiere instrumentación distinta.

\textbf{Repaso:} \hyperref[sec:r2]{Administración de redes}. \textbf{Fuentes:} \cite{rfc2131,rfc1034,rfc9293,networks}. \hyperref[q:R-029]{Volver al reactivo}.\par\medskip

\subsection*{R-030 \quad Respuesta C}\label{s:R-030}
\textbf{Razonamiento.} La autoridad corresponde a información publicada de una zona DNS.

\textbf{A) Distractor. }La temperatura no define registros.

\textbf{B) Distractor. }DNS no es un repositorio general de contraseñas.

\textbf{Repaso:} \hyperref[sec:r2]{Administración de redes}. \textbf{Fuentes:} \cite{rfc2131,rfc1034,rfc9293,networks}. \hyperref[q:R-030]{Volver al reactivo}.\par\medskip

\subsection*{R-031 \quad Respuesta A}\label{s:R-031}
\textbf{Razonamiento.} SNMP permite acceder a objetos de administración según permisos y versión.

\textbf{B) Distractor. }Coexiste con los protocolos que transportan tráfico.

\textbf{C) Distractor. }No es su función técnica.

\textbf{Repaso:} \hyperref[sec:r2]{Administración de redes}. \textbf{Fuentes:} \cite{rfc2131,rfc1034,rfc9293,networks}. \hyperref[q:R-031]{Volver al reactivo}.\par\medskip

\subsection*{R-032 \quad Respuesta B}\label{s:R-032}
\textbf{Razonamiento.} NTP estima y corrige diferencias de tiempo entre equipos.

\textbf{A) Distractor. }La segmentación de capa 2 requiere otros mecanismos.

\textbf{C) Distractor. }El cifrado de almacenamiento es otra función.

\textbf{Repaso:} \hyperref[sec:r2]{Administración de redes}. \textbf{Fuentes:} \cite{rfc2131,rfc1034,rfc9293,networks}. \hyperref[q:R-032]{Volver al reactivo}.\par\medskip

\subsection*{R-033 \quad Respuesta A}\label{s:R-033}
\textbf{Razonamiento.} Permite comparar desviaciones de tráfico, latencia y recursos.

\textbf{B) Distractor. }Una credencial no es una referencia de rendimiento.

\textbf{C) Distractor. }Las mediciones históricas no aseguran el futuro.

\textbf{Repaso:} \hyperref[sec:r2]{Administración de redes}. \textbf{Fuentes:} \cite{rfc2131,rfc1034,rfc9293,networks}. \hyperref[q:R-033]{Volver al reactivo}.\par\medskip

\subsection*{R-034 \quad Respuesta B}\label{s:R-034}
\textbf{Razonamiento.} Los puertos permiten demultiplexar conexiones hacia servicios.

\textbf{A) Distractor. }El fabricante no identifica el proceso receptor.

\textbf{C) Distractor. }La etiqueta no está en el encabezado TCP.

\textbf{Repaso:} \hyperref[sec:r2]{Administración de redes}. \textbf{Fuentes:} \cite{rfc2131,rfc1034,rfc9293,networks}. \hyperref[q:R-034]{Volver al reactivo}.\par\medskip

\subsection*{R-035 \quad Respuesta C}\label{s:R-035}
\textbf{Razonamiento.} El cliente debe renovar o dejar de usar la asignación según la concesión.

\textbf{A) Distractor. }Puede expirar y no implica titularidad permanente.

\textbf{B) Distractor. }Las claves criptográficas pertenecen a otro mecanismo.

\textbf{Repaso:} \hyperref[sec:r2]{Administración de redes}. \textbf{Fuentes:} \cite{rfc2131,rfc1034,rfc9293,networks}. \hyperref[q:R-035]{Volver al reactivo}.\par\medskip

\subsection*{R-036 \quad Respuesta B}\label{s:R-036}
\textbf{Razonamiento.} AAAA publica una dirección IPv6 asociada al nombre.

\textbf{A) Distractor. }MX identifica intercambiadores de correo.

\textbf{C) Distractor. }A es el registro usual para IPv4.

\textbf{Repaso:} \hyperref[sec:r2]{Administración de redes}. \textbf{Fuentes:} \cite{rfc2131,rfc1034,rfc9293,networks}. \hyperref[q:R-036]{Volver al reactivo}.\par\medskip

\subsection*{R-037 \quad Respuesta C}\label{s:R-037}
\textbf{Razonamiento.} Conservar configuraciones conocidas permite recuperación y auditoría.

\textbf{A) Distractor. }No modifica las capacidades del hardware.

\textbf{B) Distractor. }La copia debe validarse y probarse para ser útil.

\textbf{Repaso:} \hyperref[sec:r2]{Administración de redes}. \textbf{Fuentes:} \cite{rfc2131,rfc1034,rfc9293,networks}. \hyperref[q:R-037]{Volver al reactivo}.\par\medskip

\subsection*{R-038 \quad Respuesta A}\label{s:R-038}
\textbf{Razonamiento.} El cliente descubre, recibe oferta, solicita y obtiene confirmación.

\textbf{B) Distractor. }El servidor no ofrece antes del descubrimiento en ese flujo.

\textbf{C) Distractor. }Inicia con una solicitud que presupone información de oferta.

\textbf{Repaso:} \hyperref[sec:r2]{Administración de redes}. \textbf{Fuentes:} \cite{rfc2131,rfc1034,rfc9293,networks}. \hyperref[q:R-038]{Volver al reactivo}.\par\medskip

\subsection*{R-039 \quad Respuesta C}\label{s:R-039}
\textbf{Razonamiento.} La conectividad IP funciona; el síntoma diferencia la resolución de nombres.

\textbf{A) Distractor. }No está sustentado por el acceso exitoso mediante IP.

\textbf{B) Distractor. }Es una acción amplia sin aislar la causa.

\textbf{Repaso:} \hyperref[sec:r2]{Administración de redes}. \textbf{Fuentes:} \cite{rfc2131,rfc1034,rfc9293,networks}. \hyperref[q:R-039]{Volver al reactivo}.\par\medskip

\subsection*{R-040 \quad Respuesta A}\label{s:R-040}
\textbf{Razonamiento.} La dirección IPv4 link-local puede aparecer tras fallar la configuración DHCP.

\textbf{B) Distractor. }Ese bloque no es de direccionamiento público ordinario.

\textbf{C) Distractor. }Un certificado no explica la dirección local adquirida.

\textbf{Repaso:} \hyperref[sec:r2]{Administración de redes}. \textbf{Fuentes:} \cite{rfc2131,rfc1034,rfc9293,networks}. \hyperref[q:R-040]{Volver al reactivo}.\par\medskip

\subsection*{R-041 \quad Respuesta B}\label{s:R-041}
\textbf{Razonamiento.} Los CRC apuntan a corrupción de tramas, compatible con defectos físicos.

\textbf{A) Distractor. }El nombre no altera la integridad de bits del enlace.

\textbf{C) Distractor. }SQL opera a otro nivel y no produce ese contador físico.

\textbf{Repaso:} \hyperref[sec:r2]{Administración de redes}. \textbf{Fuentes:} \cite{rfc2131,rfc1034,rfc9293,networks}. \hyperref[q:R-041]{Volver al reactivo}.\par\medskip

\subsection*{R-042 \quad Respuesta A}\label{s:R-042}
\textbf{Razonamiento.} 250/1000=0.25.

\textbf{B) Distractor. }Es la fracción nominal restante.

\textbf{C) Distractor. }Invierte o escala incorrectamente la relación.

\textbf{Repaso:} \hyperref[sec:r2]{Administración de redes}. \textbf{Fuentes:} \cite{rfc2131,rfc1034,rfc9293,networks}. \hyperref[q:R-042]{Volver al reactivo}.\par\medskip

\subsection*{R-043 \quad Respuesta B}\label{s:R-043}
\textbf{Razonamiento.} A=990/1000×100=99\%.

\textbf{A) Distractor. }Correspondería a 900 horas.

\textbf{C) Distractor. }Correspondería a 999 horas.

\textbf{Repaso:} \hyperref[sec:r2]{Administración de redes}. \textbf{Fuentes:} \cite{rfc2131,rfc1034,rfc9293,networks}. \hyperref[q:R-043]{Volver al reactivo}.\par\medskip

\subsection*{R-044 \quad Respuesta C}\label{s:R-044}
\textbf{Razonamiento.} Las políticas pueden permitir un servicio y descartar las pruebas ICMP.

\textbf{A) Distractor. }El funcionamiento de HTTPS contradice esa conclusión.

\textbf{B) Distractor. }HTTPS se transporta sobre la pila de red.

\textbf{Repaso:} \hyperref[sec:r2]{Administración de redes}. \textbf{Fuentes:} \cite{rfc2131,rfc1034,rfc9293,networks}. \hyperref[q:R-044]{Volver al reactivo}.\par\medskip

\subsection*{R-045 \quad Respuesta B}\label{s:R-045}
\textbf{Razonamiento.} Los resolvedores pueden conservar la respuesta hasta expirar su TTL.

\textbf{A) Distractor. }Puede ser un comportamiento normal de caché.

\textbf{C) Distractor. }El DNS no exige cambiar las direcciones de enlace.

\textbf{Repaso:} \hyperref[sec:r2]{Administración de redes}. \textbf{Fuentes:} \cite{rfc2131,rfc1034,rfc9293,networks}. \hyperref[q:R-045]{Volver al reactivo}.\par\medskip

\subsection*{R-046 \quad Respuesta C}\label{s:R-046}
\textbf{Razonamiento.} SSH protege la sesión y permite autenticación segura según configuración.

\textbf{A) Distractor. }Telnet no ofrece confidencialidad de sesión.

\textbf{B) Distractor. }Expone credenciales en tránsito.

\textbf{Repaso:} \hyperref[sec:r2]{Administración de redes}. \textbf{Fuentes:} \cite{rfc2131,rfc1034,rfc9293,networks}. \hyperref[q:R-046]{Volver al reactivo}.\par\medskip

\subsection*{R-047 \quad Respuesta A}\label{s:R-047}
\textbf{Razonamiento.} La preparación permite comparar y restaurar si la modificación falla.

\textbf{B) Distractor. }Elimina la referencia necesaria para recuperar el servicio.

\textbf{C) Distractor. }Amplía el impacto e impide aislar resultados.

\textbf{Repaso:} \hyperref[sec:r2]{Administración de redes}. \textbf{Fuentes:} \cite{rfc2131,rfc1034,rfc9293,networks}. \hyperref[q:R-047]{Volver al reactivo}.\par\medskip

\subsection*{R-048 \quad Respuesta C}\label{s:R-048}
\textbf{Razonamiento.} TCP responde a pérdidas y demora; esas métricas ayudan a localizar el límite.

\textbf{A) Distractor. }No caracteriza el trayecto de red.

\textbf{B) Distractor. }No determina la tasa de retransmisiones.

\textbf{Repaso:} \hyperref[sec:r2]{Administración de redes}. \textbf{Fuentes:} \cite{rfc2131,rfc1034,rfc9293,networks}. \hyperref[q:R-048]{Volver al reactivo}.\par\medskip

\subsection*{R-049 \quad Respuesta A}\label{s:R-049}
\textbf{Razonamiento.} El siguiente salto debe poder resolverse y alcanzarse mediante el enlace local previsto.

\textbf{B) Distractor. }No arregla el primer salto IP.

\textbf{C) Distractor. }No interviene en el encaminamiento.

\textbf{Repaso:} \hyperref[sec:r2]{Administración de redes}. \textbf{Fuentes:} \cite{rfc2131,rfc1034,rfc9293,networks}. \hyperref[q:R-049]{Volver al reactivo}.\par\medskip

\subsection*{R-050 \quad Respuesta B}\label{s:R-050}
\textbf{Razonamiento.} Los servicios resuelven nombres, sincronizan relojes y centralizan mensajes respectivamente.

\textbf{A) Distractor. }No corresponde a sus finalidades.

\textbf{C) Distractor. }Intercambia las funciones de administración.

\textbf{Repaso:} \hyperref[sec:r2]{Administración de redes}. \textbf{Fuentes:} \cite{rfc2131,rfc1034,rfc9293,networks}. \hyperref[q:R-050]{Volver al reactivo}.\par\medskip

\subsection*{R-051 \quad Respuesta A}\label{s:R-051}
\textbf{Razonamiento.} La redundancia de capa 2 sin prevención de bucles puede replicar tramas.

\textbf{B) Distractor. }No controla el reenvío Ethernet.

\textbf{C) Distractor. }No explica las tormentas de difusión.

\textbf{Repaso:} \hyperref[sec:r2]{Administración de redes}. \textbf{Fuentes:} \cite{rfc2131,rfc1034,rfc9293,networks}. \hyperref[q:R-051]{Volver al reactivo}.\par\medskip

\subsection*{R-052 \quad Respuesta B}\label{s:R-052}
\textbf{Razonamiento.} Las políticas o la asimetría de retorno pueden impedir el flujo entre segmentos.

\textbf{A) Distractor. }Reduce evidencia y no resuelve la causa.

\textbf{C) Distractor. }El síntoma distingue alcance de red, no almacenamiento.

\textbf{Repaso:} \hyperref[sec:r2]{Administración de redes}. \textbf{Fuentes:} \cite{rfc2131,rfc1034,rfc9293,networks}. \hyperref[q:R-052]{Volver al reactivo}.\par\medskip

\subsection*{R-053 \quad Respuesta C}\label{s:R-053}
\textbf{Razonamiento.} El diagnóstico debe acotar exposición y proteger evidencia.

\textbf{A) Distractor. }Expone información sin necesidad.

\textbf{B) Distractor. }Amplía el riesgo y no responde a un objetivo definido.

\textbf{Repaso:} \hyperref[sec:r2]{Administración de redes}. \textbf{Fuentes:} \cite{rfc2131,rfc1034,rfc9293,networks}. \hyperref[q:R-053]{Volver al reactivo}.\par\medskip

\subsection*{R-054 \quad Respuesta B}\label{s:R-054}
\textbf{Razonamiento.} La distribución y las ventanas cortas permiten ver episodios ocultos por la media.

\textbf{A) Distractor. }Descarta justamente el fenómeno a investigar.

\textbf{C) Distractor. }Oculta aún más los picos.

\textbf{Repaso:} \hyperref[sec:r2]{Administración de redes}. \textbf{Fuentes:} \cite{rfc2131,rfc1034,rfc9293,networks}. \hyperref[q:R-054]{Volver al reactivo}.\par\medskip

\subsection*{R-055 \quad Respuesta C}\label{s:R-055}
\textbf{Razonamiento.} La potencia no describe por sí sola ruido, retransmisiones ni competencia por el medio.

\textbf{A) Distractor. }Ya es buena y no agota las causas.

\textbf{B) Distractor. }No determina la calidad radioeléctrica.

\textbf{Repaso:} \hyperref[sec:r2]{Administración de redes}. \textbf{Fuentes:} \cite{rfc2131,rfc1034,rfc9293,networks}. \hyperref[q:R-055]{Volver al reactivo}.\par\medskip

\subsection*{R-056 \quad Respuesta A}\label{s:R-056}
\textbf{Razonamiento.} La prueba y el despliegue acotado limitan el impacto de incompatibilidades.

\textbf{B) Distractor. }Multiplica el alcance de un fallo.

\textbf{C) Distractor. }Puede introducir problemas previsibles de soporte.

\textbf{Repaso:} \hyperref[sec:r2]{Administración de redes}. \textbf{Fuentes:} \cite{rfc2131,rfc1034,rfc9293,networks}. \hyperref[q:R-056]{Volver al reactivo}.\par\medskip

\subsection*{R-057 \quad Respuesta C}\label{s:R-057}
\textbf{Razonamiento.} La saturación reduce margen para ráfagas y puede acumular espera.

\textbf{A) Distractor. }Las colas aumentan demora.

\textbf{B) Distractor. }La ocupación sola no determina intención del tráfico.

\textbf{Repaso:} \hyperref[sec:r2]{Administración de redes}. \textbf{Fuentes:} \cite{rfc2131,rfc1034,rfc9293,networks}. \hyperref[q:R-057]{Volver al reactivo}.\par\medskip

\subsection*{R-058 \quad Respuesta A}\label{s:R-058}
\textbf{Razonamiento.} Una vía fuera de banda o reversión programada evita quedar sin acceso.

\textbf{B) Distractor. }La sesión puede perderse y no restablecerse.

\textbf{C) Distractor. }Destruye medios de recuperación.

\textbf{Repaso:} \hyperref[sec:r2]{Administración de redes}. \textbf{Fuentes:} \cite{rfc2131,rfc1034,rfc9293,networks}. \hyperref[q:R-058]{Volver al reactivo}.\par\medskip

\subsection*{R-059 \quad Respuesta B}\label{s:R-059}
\textbf{Razonamiento.} Relojes coherentes y una referencia horaria explícita permiten ordenar evidencias.

\textbf{A) Distractor. }Eliminar las marcas destruye información útil.

\textbf{C) Distractor. }Los monitores no alinean las marcas de tiempo.

\textbf{Repaso:} \hyperref[sec:r2]{Administración de redes}. \textbf{Fuentes:} \cite{rfc2131,rfc1034,rfc9293,networks}. \hyperref[q:R-059]{Volver al reactivo}.\par\medskip

\section{Seguridad y evaluación en redes}

\subsection*{R-060 \quad Respuesta B}\label{s:R-060}
\textbf{Razonamiento.} La confidencialidad restringe quién conoce la información.

\textbf{A) Distractor. }La latencia es una medida temporal.

\textbf{C) Distractor. }La disponibilidad permite acceder al servicio cuando se necesita.

\textbf{Repaso:} \hyperref[sec:r3]{Seguridad y evaluación en redes}. \textbf{Fuentes:} \cite{nistsecurity,owasp,aes,tls}. \hyperref[q:R-060]{Volver al reactivo}.\par\medskip

\subsection*{R-061 \quad Respuesta C}\label{s:R-061}
\textbf{Razonamiento.} La integridad se refiere a mantener y verificar datos sin modificaciones indebidas.

\textbf{A) Distractor. }La escalabilidad describe crecimiento de capacidad.

\textbf{B) Distractor. }La tasa de transmisión no demuestra integridad.

\textbf{Repaso:} \hyperref[sec:r3]{Seguridad y evaluación en redes}. \textbf{Fuentes:} \cite{nistsecurity,owasp,aes,tls}. \hyperref[q:R-061]{Volver al reactivo}.\par\medskip

\subsection*{R-062 \quad Respuesta B}\label{s:R-062}
\textbf{Razonamiento.} La autenticación prueba una identidad mediante evidencia apropiada.

\textbf{A) Distractor. }La identidad no implica infalibilidad.

\textbf{C) Distractor. }La autorización decide permisos y puede ser restrictiva.

\textbf{Repaso:} \hyperref[sec:r3]{Seguridad y evaluación en redes}. \textbf{Fuentes:} \cite{nistsecurity,owasp,aes,tls}. \hyperref[q:R-062]{Volver al reactivo}.\par\medskip

\subsection*{R-063 \quad Respuesta C}\label{s:R-063}
\textbf{Razonamiento.} Reduce exposición y limita el impacto de una cuenta comprometida.

\textbf{A) Distractor. }Entrega más derechos de los necesarios.

\textbf{B) Distractor. }Elimina atribución y concentra privilegios.

\textbf{Repaso:} \hyperref[sec:r3]{Seguridad y evaluación en redes}. \textbf{Fuentes:} \cite{nistsecurity,owasp,aes,tls}. \hyperref[q:R-063]{Volver al reactivo}.\par\medskip

\subsection*{R-064 \quad Respuesta A}\label{s:R-064}
\textbf{Razonamiento.} AES usa una clave secreta compartida y bloques de 128 bits.

\textbf{B) Distractor. }Las firmas requieren otros algoritmos y un par de claves.

\textbf{C) Distractor. }No está diseñado como compresor.

\textbf{Repaso:} \hyperref[sec:r3]{Seguridad y evaluación en redes}. \textbf{Fuentes:} \cite{aes}. \hyperref[q:R-064]{Volver al reactivo}.\par\medskip

\subsection*{R-065 \quad Respuesta C}\label{s:R-065}
\textbf{Razonamiento.} No existe una operación de descifrado del resumen que reconstruya arbitrariamente el mensaje.

\textbf{A) Distractor. }La longitud del resumen suele ser fija.

\textbf{B) Distractor. }El resumen no funciona como un texto cifrado reversible.

\textbf{Repaso:} \hyperref[sec:r3]{Seguridad y evaluación en redes}. \textbf{Fuentes:} \cite{nistsecurity,owasp,aes,tls}. \hyperref[q:R-065]{Volver al reactivo}.\par\medskip

\subsection*{R-066 \quad Respuesta A}\label{s:R-066}
\textbf{Razonamiento.} Puede combinar conocimiento con posesión o inherencia.

\textbf{B) Distractor. }Dos secretos memorizados siguen siendo del mismo factor.

\textbf{C) Distractor. }La reutilización no añade un factor.

\textbf{Repaso:} \hyperref[sec:r3]{Seguridad y evaluación en redes}. \textbf{Fuentes:} \cite{nistsecurity,owasp,aes,tls}. \hyperref[q:R-066]{Volver al reactivo}.\par\medskip

\subsection*{R-067 \quad Respuesta B}\label{s:R-067}
\textbf{Razonamiento.} Filtra o inspecciona comunicaciones con criterios definidos.

\textbf{A) Distractor. }No corrige vulnerabilidades internas del software.

\textbf{C) Distractor. }No elimina todos los ataques de ingeniería social.

\textbf{Repaso:} \hyperref[sec:r3]{Seguridad y evaluación en redes}. \textbf{Fuentes:} \cite{nistsecurity,owasp,aes,tls}. \hyperref[q:R-067]{Volver al reactivo}.\par\medskip

\subsection*{R-068 \quad Respuesta A}\label{s:R-068}
\textbf{Razonamiento.} El atacante suplanta una entidad o contexto para manipular a la víctima.

\textbf{B) Distractor. }El reparto de carga no define el engaño.

\textbf{C) Distractor. }No se refiere al mantenimiento físico.

\textbf{Repaso:} \hyperref[sec:r3]{Seguridad y evaluación en redes}. \textbf{Fuentes:} \cite{nistsecurity,owasp,aes,tls}. \hyperref[q:R-068]{Volver al reactivo}.\par\medskip

\subsection*{R-069 \quad Respuesta B}\label{s:R-069}
\textbf{Razonamiento.} La copia reduce el impacto si su restauración es viable y probada.

\textbf{A) Distractor. }No controla quién entra al sistema.

\textbf{C) Distractor. }No bloquea necesariamente el acceso del atacante.

\textbf{Repaso:} \hyperref[sec:r3]{Seguridad y evaluación en redes}. \textbf{Fuentes:} \cite{nistsecurity,owasp,aes,tls}. \hyperref[q:R-069]{Volver al reactivo}.\par\medskip

\subsection*{R-070 \quad Respuesta C}\label{s:R-070}
\textbf{Razonamiento.} La combinación reduce dependencia de un único punto de protección.

\textbf{A) Distractor. }Una sola barrera puede fallar.

\textbf{B) Distractor. }La protección perimetral no vuelve innecesarios otros controles.

\textbf{Repaso:} \hyperref[sec:r3]{Seguridad y evaluación en redes}. \textbf{Fuentes:} \cite{nistsecurity,owasp,aes,tls}. \hyperref[q:R-070]{Volver al reactivo}.\par\medskip

\subsection*{R-071 \quad Respuesta B}\label{s:R-071}
\textbf{Razonamiento.} El riesgo surge al relacionar agentes o eventos con debilidades e impacto.

\textbf{A) Distractor. }La debilidad puede existir antes de ser explotada.

\textbf{C) Distractor. }Desempeñan papeles diferentes en el análisis.

\textbf{Repaso:} \hyperref[sec:r3]{Seguridad y evaluación en redes}. \textbf{Fuentes:} \cite{nistsecurity,owasp,aes,tls}. \hyperref[q:R-071]{Volver al reactivo}.\par\medskip

\subsection*{R-072 \quad Respuesta C}\label{s:R-072}
\textbf{Razonamiento.} Se concede la capacidad mínima requerida por la tarea.

\textbf{A) Distractor. }Permite acciones ajenas a consultar reportes.

\textbf{B) Distractor. }Incluye modificaciones no solicitadas.

\textbf{Repaso:} \hyperref[sec:r3]{Seguridad y evaluación en redes}. \textbf{Fuentes:} \cite{nistsecurity,owasp,aes,tls}. \hyperref[q:R-072]{Volver al reactivo}.\par\medskip

\subsection*{R-073 \quad Respuesta A}\label{s:R-073}
\textbf{Razonamiento.} Los parámetros separan datos de la estructura de la consulta.

\textbf{B) Distractor. }No cambia la interpretación de SQL.

\textbf{C) Distractor. }El tamaño del nombre no evita la concatenación peligrosa.

\textbf{Repaso:} \hyperref[sec:r3]{Seguridad y evaluación en redes}. \textbf{Fuentes:} \cite{owasp,postgres}. \hyperref[q:R-073]{Volver al reactivo}.\par\medskip

\subsection*{R-074 \quad Respuesta C}\label{s:R-074}
\textbf{Razonamiento.} El canal cifrado necesita autenticar también al extremo previsto.

\textbf{A) Distractor. }El cifrado sin identidad puede conectar con un impostor.

\textbf{B) Distractor. }Elimina la protección de autenticación del servidor.

\textbf{Repaso:} \hyperref[sec:r3]{Seguridad y evaluación en redes}. \textbf{Fuentes:} \cite{tls}. \hyperref[q:R-074]{Volver al reactivo}.\par\medskip

\subsection*{R-075 \quad Respuesta A}\label{s:R-075}
\textbf{Razonamiento.} El costo de derivación y la sal dificultan ataques masivos a las credenciales almacenadas.

\textbf{B) Distractor. }Si se accede a la tabla se revelan directamente.

\textbf{C) Distractor. }Facilita comparación y pruebas rápidas de diccionario.

\textbf{Repaso:} \hyperref[sec:r3]{Seguridad y evaluación en redes}. \textbf{Fuentes:} \cite{nistsecurity,owasp,aes,tls}. \hyperref[q:R-075]{Volver al reactivo}.\par\medskip

\subsection*{R-076 \quad Respuesta B}\label{s:R-076}
\textbf{Razonamiento.} El código de autenticación depende del mensaje y la clave compartida.

\textbf{A) Distractor. }Base64 representa datos, sin aportar autenticación.

\textbf{C) Distractor. }Cualquiera puede recalcular un hash público tras modificar el texto.

\textbf{Repaso:} \hyperref[sec:r3]{Seguridad y evaluación en redes}. \textbf{Fuentes:} \cite{nistsecurity,owasp,aes,tls}. \hyperref[q:R-076]{Volver al reactivo}.\par\medskip

\subsection*{R-077 \quad Respuesta A}\label{s:R-077}
\textbf{Razonamiento.} La clave pública verifica una firma generada con la privada correspondiente.

\textbf{B) Distractor. }Cualquier poseedor del secreto compartido podría producir el mismo tipo de evidencia.

\textbf{C) Distractor. }Comprimir no identifica al firmante.

\textbf{Repaso:} \hyperref[sec:r3]{Seguridad y evaluación en redes}. \textbf{Fuentes:} \cite{nistsecurity,owasp,aes,tls}. \hyperref[q:R-077]{Volver al reactivo}.\par\medskip

\subsection*{R-078 \quad Respuesta B}\label{s:R-078}
\textbf{Razonamiento.} La separación lógica requiere una política efectiva de tráfico entre segmentos.

\textbf{A) Distractor. }Facilita precisamente el acceso no autorizado.

\textbf{C) Distractor. }El SSID por sí solo no demuestra aislamiento.

\textbf{Repaso:} \hyperref[sec:r3]{Seguridad y evaluación en redes}. \textbf{Fuentes:} \cite{nistsecurity,owasp,aes,tls}. \hyperref[q:R-078]{Volver al reactivo}.\par\medskip

\subsection*{R-079 \quad Respuesta C}\label{s:R-079}
\textbf{Razonamiento.} La verificación y el tratamiento proporcionado reducen riesgo sin cambios ciegos.

\textbf{A) Distractor. }Puede causar daño y exceder el alcance permitido.

\textbf{B) Distractor. }La ausencia de incidente conocido no elimina la debilidad.

\textbf{Repaso:} \hyperref[sec:r3]{Seguridad y evaluación en redes}. \textbf{Fuentes:} \cite{nistsecurity,owasp,aes,tls}. \hyperref[q:R-079]{Volver al reactivo}.\par\medskip

\subsection*{R-080 \quad Respuesta B}\label{s:R-080}
\textbf{Razonamiento.} La revocación limita uso indebido y la verificación evita toma de cuenta.

\textbf{A) Distractor. }Rompe identidad y trazabilidad.

\textbf{C) Distractor. }Amplía innecesariamente la exposición global.

\textbf{Repaso:} \hyperref[sec:r3]{Seguridad y evaluación en redes}. \textbf{Fuentes:} \cite{nistsecurity,owasp,aes,tls}. \hyperref[q:R-080]{Volver al reactivo}.\par\medskip

\subsection*{R-081 \quad Respuesta C}\label{s:R-081}
\textbf{Razonamiento.} La separación reduce la posibilidad de que el incidente destruya todas las versiones.

\textbf{A) Distractor. }Comparte el mismo dominio de compromiso.

\textbf{B) Distractor. }Un acceso directo no es una copia independiente.

\textbf{Repaso:} \hyperref[sec:r3]{Seguridad y evaluación en redes}. \textbf{Fuentes:} \cite{nistsecurity,owasp,aes,tls}. \hyperref[q:R-081]{Volver al reactivo}.\par\medskip

\subsection*{R-082 \quad Respuesta A}\label{s:R-082}
\textbf{Razonamiento.} La prevención incorpora una capacidad de intervención sobre el tráfico.

\textbf{B) Distractor. }No describe detección ni prevención de intrusiones.

\textbf{C) Distractor. }Los sistemas pueden producir errores de clasificación.

\textbf{Repaso:} \hyperref[sec:r3]{Seguridad y evaluación en redes}. \textbf{Fuentes:} \cite{nistsecurity,owasp,aes,tls}. \hyperref[q:R-082]{Volver al reactivo}.\par\medskip

\subsection*{R-083 \quad Respuesta C}\label{s:R-083}
\textbf{Razonamiento.} La correlación permite distinguir abuso y activar controles apropiados.

\textbf{A) Distractor. }Destruye evidencia relevante.

\textbf{B) Distractor. }Expone las credenciales y no es una prueba segura.

\textbf{Repaso:} \hyperref[sec:r3]{Seguridad y evaluación en redes}. \textbf{Fuentes:} \cite{nistsecurity,owasp,aes,tls}. \hyperref[q:R-083]{Volver al reactivo}.\par\medskip

\subsection*{R-084 \quad Respuesta A}\label{s:R-084}
\textbf{Razonamiento.} Primero se comprende el evento, después se limita y elimina la causa y finalmente se restaura.

\textbf{B) Distractor. }La recuperación requiere conocer y tratar el incidente.

\textbf{C) Distractor. }No se contiene de forma informada sin detección ni se omite análisis.

\textbf{Repaso:} \hyperref[sec:r3]{Seguridad y evaluación en redes}. \textbf{Fuentes:} \cite{nistsecurity,owasp,aes,tls}. \hyperref[q:R-084]{Volver al reactivo}.\par\medskip

\subsection*{R-085 \quad Respuesta B}\label{s:R-085}
\textbf{Razonamiento.} Las tres propiedades protegen divulgación, alteración y acceso oportuno respectivamente.

\textbf{A) Distractor. }No coincide con sus definiciones.

\textbf{C) Distractor. }Intercambia los objetivos de la tríada.

\textbf{Repaso:} \hyperref[sec:r3]{Seguridad y evaluación en redes}. \textbf{Fuentes:} \cite{nistsecurity,owasp,aes,tls}. \hyperref[q:R-085]{Volver al reactivo}.\par\medskip

\subsection*{R-086 \quad Respuesta A}\label{s:R-086}
\textbf{Razonamiento.} El alcance define activos, técnicas y límites permitidos.

\textbf{B) Distractor. }La accesibilidad no equivale a autorización.

\textbf{C) Distractor. }La trazabilidad es necesaria para controlar y explicar resultados.

\textbf{Repaso:} \hyperref[sec:r3]{Seguridad y evaluación en redes}. \textbf{Fuentes:} \cite{nistsecurity,owasp,aes,tls}. \hyperref[q:R-086]{Volver al reactivo}.\par\medskip

\subsection*{R-087 \quad Respuesta B}\label{s:R-087}
\textbf{Razonamiento.} Reduce quién puede alcanzar la interfaz sensible y protege la sesión.

\textbf{A) Distractor. }Amplía la superficie expuesta.

\textbf{C) Distractor. }Ocultar información no sustituye control de acceso.

\textbf{Repaso:} \hyperref[sec:r3]{Seguridad y evaluación en redes}. \textbf{Fuentes:} \cite{nistsecurity,owasp,aes,tls}. \hyperref[q:R-087]{Volver al reactivo}.\par\medskip

\subsection*{R-088 \quad Respuesta C}\label{s:R-088}
\textbf{Razonamiento.} Limitar comunicaciones reduce oportunidades de abuso y exfiltración.

\textbf{A) Distractor. }Habilita tráfico no requerido.

\textbf{B) Distractor. }Otorga privilegios ajenos a su función.

\textbf{Repaso:} \hyperref[sec:r3]{Seguridad y evaluación en redes}. \textbf{Fuentes:} \cite{nistsecurity,owasp,aes,tls}. \hyperref[q:R-088]{Volver al reactivo}.\par\medskip

\subsection*{R-089 \quad Respuesta B}\label{s:R-089}
\textbf{Razonamiento.} La relevancia depende de exposición, posibilidad de materialización y consecuencias.

\textbf{A) Distractor. }La marca no sustituye el análisis del entorno.

\textbf{C) Distractor. }La longitud textual no mide riesgo.

\textbf{Repaso:} \hyperref[sec:r3]{Seguridad y evaluación en redes}. \textbf{Fuentes:} \cite{nistsecurity,owasp,aes,tls}. \hyperref[q:R-089]{Volver al reactivo}.\par\medskip

\subsection*{R-090 \quad Respuesta C}\label{s:R-090}
\textbf{Razonamiento.} La clave comprometida ya no sirve como secreto de autenticación.

\textbf{A) Distractor. }No afecta la confianza criptográfica.

\textbf{B) Distractor. }La fecha de expiración no revierte la exposición.

\textbf{Repaso:} \hyperref[sec:r3]{Seguridad y evaluación en redes}. \textbf{Fuentes:} \cite{tls,nistsecurity}. \hyperref[q:R-090]{Volver al reactivo}.\par\medskip

\subsection*{R-091 \quad Respuesta A}\label{s:R-091}
\textbf{Razonamiento.} El éxito de escritura no prueba integridad funcional ni tiempo de recuperación.

\textbf{B) Distractor. }La falta de prueba no demuestra un ataque.

\textbf{C) Distractor. }La recuperación y el acceso son controles distintos.

\textbf{Repaso:} \hyperref[sec:r3]{Seguridad y evaluación en redes}. \textbf{Fuentes:} \cite{nistsecurity,owasp,aes,tls}. \hyperref[q:R-091]{Volver al reactivo}.\par\medskip

\subsection*{R-092 \quad Respuesta C}\label{s:R-092}
\textbf{Razonamiento.} La contención reduce exposición y mantiene datos para comprender el incidente.

\textbf{A) Distractor. }Imposibilita reconstruir acciones.

\textbf{B) Distractor. }Aumenta las oportunidades de propagación.

\textbf{Repaso:} \hyperref[sec:r3]{Seguridad y evaluación en redes}. \textbf{Fuentes:} \cite{nistsecurity,owasp,aes,tls}. \hyperref[q:R-092]{Volver al reactivo}.\par\medskip

\subsection*{R-093 \quad Respuesta A}\label{s:R-093}
\textbf{Razonamiento.} La identidad compartida dificulta saber quién realizó cada operación.

\textbf{B) Distractor. }No modifica la tasa física del enlace.

\textbf{C) Distractor. }La longitud de direcciones no depende de cuentas.

\textbf{Repaso:} \hyperref[sec:r3]{Seguridad y evaluación en redes}. \textbf{Fuentes:} \cite{nistsecurity,owasp,aes,tls}. \hyperref[q:R-093]{Volver al reactivo}.\par\medskip

\subsection*{R-094 \quad Respuesta B}\label{s:R-094}
\textbf{Razonamiento.} La protección del canal no corrige un extremo comprometido.

\textbf{A) Distractor. }Un atacante en el extremo puede generar tráfico dentro del túnel.

\textbf{C) Distractor. }El cifrado de tránsito no desinfecta el sistema.

\textbf{Repaso:} \hyperref[sec:r3]{Seguridad y evaluación en redes}. \textbf{Fuentes:} \cite{nistsecurity,owasp,aes,tls}. \hyperref[q:R-094]{Volver al reactivo}.\par\medskip

\subsection*{R-095 \quad Respuesta A}\label{s:R-095}
\textbf{Razonamiento.} La firma solo protege si el receptor la verifica con una clave confiable.

\textbf{B) Distractor. }La compresión no prueba origen.

\textbf{C) Distractor. }El nombre no valida el contenido.

\textbf{Repaso:} \hyperref[sec:r3]{Seguridad y evaluación en redes}. \textbf{Fuentes:} \cite{nistsecurity,owasp,aes,tls}. \hyperref[q:R-095]{Volver al reactivo}.\par\medskip

\subsection*{R-096 \quad Respuesta B}\label{s:R-096}
\textbf{Razonamiento.} RPO se refiere al punto de recuperación; RTO al tiempo objetivo para restablecer.

\textbf{A) Distractor. }Los objetivos no prescriben ese número de copias.

\textbf{C) Distractor. }Intercambia las dos métricas.

\textbf{Repaso:} \hyperref[sec:r3]{Seguridad y evaluación en redes}. \textbf{Fuentes:} \cite{nistsecurity,owasp,aes,tls}. \hyperref[q:R-096]{Volver al reactivo}.\par\medskip

\subsection*{R-097 \quad Respuesta C}\label{s:R-097}
\textbf{Razonamiento.} Se alerta sobre una condición de ataque que realmente no existe.

\textbf{A) Distractor. }Sería no detectar un ataque real.

\textbf{B) Distractor. }Requiere que la amenaza sí esté presente.

\textbf{Repaso:} \hyperref[sec:r3]{Seguridad y evaluación en redes}. \textbf{Fuentes:} \cite{nistsecurity,owasp,aes,tls}. \hyperref[q:R-097]{Volver al reactivo}.\par\medskip

\subsection*{R-098 \quad Respuesta B}\label{s:R-098}
\textbf{Razonamiento.} La verificación matemática no prueba por sí sola de quién es la clave.

\textbf{A) Distractor. }No afecta la vinculación de identidad.

\textbf{C) Distractor. }No establece la identidad del titular.

\textbf{Repaso:} \hyperref[sec:r3]{Seguridad y evaluación en redes}. \textbf{Fuentes:} \cite{nistsecurity,owasp,aes,tls}. \hyperref[q:R-098]{Volver al reactivo}.\par\medskip

\subsection*{R-099 \quad Respuesta C}\label{s:R-099}
\textbf{Razonamiento.} La seguridad depende de primitivas analizadas y de su uso y operación correctos.

\textbf{A) Distractor. }El ocultamiento del diseño no ofrece evidencia suficiente.

\textbf{B) Distractor. }La codificación es reversible sin secreto.

\textbf{Repaso:} \hyperref[sec:r3]{Seguridad y evaluación en redes}. \textbf{Fuentes:} \cite{nistsecurity,owasp,aes,tls}. \hyperref[q:R-099]{Volver al reactivo}.\par\medskip

\subsection*{R-100 \quad Respuesta A}\label{s:R-100}
\textbf{Razonamiento.} Un control debe proteger sin impedir innecesariamente los servicios autorizados.

\textbf{B) Distractor. }Más bloqueo no equivale a una política correcta.

\textbf{C) Distractor. }Es una reacción desproporcionada que no corrige la regla específica.

\textbf{Repaso:} \hyperref[sec:r3]{Seguridad y evaluación en redes}. \textbf{Fuentes:} \cite{nistsecurity,owasp,aes,tls}. \hyperref[q:R-100]{Volver al reactivo}.\par\medskip

\section{Ingeniería de software}

\subsection*{S-001 \quad Respuesta A}\label{s:S-001}
\textbf{Razonamiento.} Especifica servicios, cálculos o respuestas observables del producto.

\textbf{B) Distractor. }No establece una función del sistema.

\textbf{C) Distractor. }El presupuesto es una restricción de gestión.

\textbf{Repaso:} \hyperref[sec:s1]{Ingeniería de software}. \textbf{Fuentes:} \cite{scrum,swebok}. \hyperref[q:S-001]{Volver al reactivo}.\par\medskip

\subsection*{S-002 \quad Respuesta B}\label{s:S-002}
\textbf{Razonamiento.} Fija una métrica, umbral y contexto que permiten una prueba.

\textbf{A) Distractor. }No especifica una propiedad comprobable del producto.

\textbf{C) Distractor. }Es subjetivo y no establece criterio de aceptación.

\textbf{Repaso:} \hyperref[sec:s1]{Ingeniería de software}. \textbf{Fuentes:} \cite{scrum,swebok}. \hyperref[q:S-002]{Volver al reactivo}.\par\medskip

\subsection*{S-003 \quad Respuesta A}\label{s:S-003}
\textbf{Razonamiento.} Permite justificar elementos y evaluar impacto de cambios.

\textbf{B) Distractor. }La falta de identidad dificulta seguir relaciones.

\textbf{C) Distractor. }Un cambio exige revisar pruebas relacionadas.

\textbf{Repaso:} \hyperref[sec:s1]{Ingeniería de software}. \textbf{Fuentes:} \cite{scrum,swebok}. \hyperref[q:S-003]{Volver al reactivo}.\par\medskip

\subsection*{S-004 \quad Respuesta B}\label{s:S-004}
\textbf{Razonamiento.} La validación pregunta si se construye el producto adecuado para el uso previsto.

\textbf{A) Distractor. }El tamaño no demuestra utilidad ni adecuación.

\textbf{C) Distractor. }Compilar comprueba aspectos limitados de construcción.

\textbf{Repaso:} \hyperref[sec:s1]{Ingeniería de software}. \textbf{Fuentes:} \cite{scrum,swebok}. \hyperref[q:S-004]{Volver al reactivo}.\par\medskip

\subsection*{S-005 \quad Respuesta C}\label{s:S-005}
\textbf{Razonamiento.} El actor describe un rol humano o sistema externo respecto del límite analizado.

\textbf{A) Distractor. }Las clases representan otra vista del diseño.

\textbf{B) Distractor. }No se restringe a hardware de procesamiento.

\textbf{Repaso:} \hyperref[sec:s1]{Ingeniería de software}. \textbf{Fuentes:} \cite{scrum,swebok}. \hyperref[q:S-005]{Volver al reactivo}.\par\medskip

\subsection*{S-006 \quad Respuesta B}\label{s:S-006}
\textbf{Razonamiento.} Un conjunto enfocado facilita entender y cambiar el módulo.

\textbf{A) Distractor. }El acoplamiento innecesario amplía el impacto de cambios.

\textbf{C) Distractor. }La mezcla reduce la claridad de propósito.

\textbf{Repaso:} \hyperref[sec:s1]{Ingeniería de software}. \textbf{Fuentes:} \cite{scrum,swebok}. \hyperref[q:S-006]{Volver al reactivo}.\par\medskip

\subsection*{S-007 \quad Respuesta C}\label{s:S-007}
\textbf{Razonamiento.} Las interfaces delimitadas permiten cambios con menor propagación.

\textbf{A) Distractor. }Un sistema puede requerir colaboración legítima.

\textbf{B) Distractor. }Duplicar introduce otros costos y no es una regla general.

\textbf{Repaso:} \hyperref[sec:s1]{Ingeniería de software}. \textbf{Fuentes:} \cite{scrum,swebok}. \hyperref[q:S-007]{Volver al reactivo}.\par\medskip

\subsection*{S-008 \quad Respuesta A}\label{s:S-008}
\textbf{Razonamiento.} El Product Backlog es una lista emergente y ordenada de lo necesario para el producto.

\textbf{B) Distractor. }No es un artefacto Scrum ni un objetivo válido.

\textbf{C) Distractor. }El backlog expresa trabajo del producto, no solo tiempo pasado.

\textbf{Repaso:} \hyperref[sec:s1]{Ingeniería de software}. \textbf{Fuentes:} \cite{scrum}. \hyperref[q:S-008]{Volver al reactivo}.\par\medskip

\subsection*{S-009 \quad Respuesta C}\label{s:S-009}
\textbf{Razonamiento.} El Product Owner es responsable de la gestión efectiva del Product Backlog y del valor.

\textbf{A) Distractor. }Los interesados aportan información pero no sustituyen esa responsabilidad.

\textbf{B) Distractor. }No es una responsabilidad del marco Scrum.

\textbf{Repaso:} \hyperref[sec:s1]{Ingeniería de software}. \textbf{Fuentes:} \cite{scrum}. \hyperref[q:S-009]{Volver al reactivo}.\par\medskip

\subsection*{S-010 \quad Respuesta A}\label{s:S-010}
\textbf{Razonamiento.} Cambian el alcance y los tipos de defecto que se buscan.

\textbf{B) Distractor. }Invierte o inventa restricciones ajenas a esos niveles.

\textbf{C) Distractor. }La colaboración entre componentes añade riesgos específicos.

\textbf{Repaso:} \hyperref[sec:s1]{Ingeniería de software}. \textbf{Fuentes:} \cite{scrum,swebok}. \hyperref[q:S-010]{Volver al reactivo}.\par\medskip

\subsection*{S-011 \quad Respuesta B}\label{s:S-011}
\textbf{Razonamiento.} Incluye los dos límites y los valores contiguos fuera del rango.

\textbf{A) Distractor. }No se concentra en los límites especificados.

\textbf{C) Distractor. }Solo prueba valores interiores.

\textbf{Repaso:} \hyperref[sec:s1]{Ingeniería de software}. \textbf{Fuentes:} \cite{scrum,swebok}. \hyperref[q:S-011]{Volver al reactivo}.\par\medskip

\subsection*{S-012 \quad Respuesta A}\label{s:S-012}
\textbf{Razonamiento.} Cada clase agrupa entradas con comportamiento esperado común.

\textbf{B) Distractor. }Puede omitir categorías completas.

\textbf{C) Distractor. }La repetición no cubre reglas distintas.

\textbf{Repaso:} \hyperref[sec:s1]{Ingeniería de software}. \textbf{Fuentes:} \cite{scrum,swebok}. \hyperref[q:S-012]{Volver al reactivo}.\par\medskip

\subsection*{S-013 \quad Respuesta B}\label{s:S-013}
\textbf{Razonamiento.} Los incrementos permiten contrastar valor y adaptar prioridades.

\textbf{A) Distractor. }Impide aprender del uso y necesidades.

\textbf{C) Distractor. }Retrasa la detección de malentendidos.

\textbf{Repaso:} \hyperref[sec:s1]{Ingeniería de software}. \textbf{Fuentes:} \cite{scrum,swebok}. \hyperref[q:S-013]{Volver al reactivo}.\par\medskip

\subsection*{S-014 \quad Respuesta C}\label{s:S-014}
\textbf{Razonamiento.} La rapidez debe convertirse en un criterio verificable en un contexto concreto.

\textbf{A) Distractor. }No permite distinguir cumplimiento de incumplimiento.

\textbf{B) Distractor. }Las líneas no equivalen a tiempo de respuesta.

\textbf{Repaso:} \hyperref[sec:s1]{Ingeniería de software}. \textbf{Fuentes:} \cite{scrum,swebok}. \hyperref[q:S-014]{Volver al reactivo}.\par\medskip

\subsection*{S-015 \quad Respuesta B}\label{s:S-015}
\textbf{Razonamiento.} El diagrama de secuencia representa colaboración temporal mediante mensajes.

\textbf{A) Distractor. }Las relaciones estáticas no describen por sí solas el orden temporal.

\textbf{C) Distractor. }El despliegue muestra nodos y artefactos físicos.

\textbf{Repaso:} \hyperref[sec:s1]{Ingeniería de software}. \textbf{Fuentes:} \cite{scrum,swebok}. \hyperref[q:S-015]{Volver al reactivo}.\par\medskip

\subsection*{S-016 \quad Respuesta C}\label{s:S-016}
\textbf{Razonamiento.} Se elimina un defecto observado en el comportamiento existente.

\textbf{A) Distractor. }Adaptativo responde a cambios del entorno.

\textbf{B) Distractor. }Preventivo busca evitar problemas futuros sin ese defecto ya reportado.

\textbf{Repaso:} \hyperref[sec:s1]{Ingeniería de software}. \textbf{Fuentes:} \cite{scrum,swebok}. \hyperref[q:S-016]{Volver al reactivo}.\par\medskip

\subsection*{S-017 \quad Respuesta A}\label{s:S-017}
\textbf{Razonamiento.} Se acomoda el sistema a un cambio externo de plataforma o dependencia.

\textbf{B) Distractor. }No implica que el código previo incumpliera su contexto anterior.

\textbf{C) Distractor. }El requisito puede seguir siendo válido.

\textbf{Repaso:} \hyperref[sec:s1]{Ingeniería de software}. \textbf{Fuentes:} \cite{scrum,swebok}. \hyperref[q:S-017]{Volver al reactivo}.\par\medskip

\subsection*{S-018 \quad Respuesta C}\label{s:S-018}
\textbf{Razonamiento.} 200×300+10000=70000.

\textbf{A) Distractor. }Omite el costo adicional.

\textbf{B) Distractor. }Multiplica o suma sin respetar los componentes del presupuesto.

\textbf{Repaso:} \hyperref[sec:s1]{Ingeniería de software}. \textbf{Fuentes:} \cite{scrum,swebok}. \hyperref[q:S-018]{Volver al reactivo}.\par\medskip

\subsection*{S-019 \quad Respuesta A}\label{s:S-019}
\textbf{Razonamiento.} La dependencia secuencial suma duraciones: 5+3=8.

\textbf{B) Distractor. }Supondría ejecución paralela que no está permitida.

\textbf{C) Distractor. }La diferencia no expresa la duración de la cadena.

\textbf{Repaso:} \hyperref[sec:s1]{Ingeniería de software}. \textbf{Fuentes:} \cite{scrum,swebok}. \hyperref[q:S-019]{Volver al reactivo}.\par\medskip

\subsection*{S-020 \quad Respuesta B}\label{s:S-020}
\textbf{Razonamiento.} Se analizan capacidades, restricciones e integración tecnológica.

\textbf{A) Distractor. }La tipografía no cubre los recursos técnicos.

\textbf{C) Distractor. }La apariencia no demuestra viabilidad de implementación.

\textbf{Repaso:} \hyperref[sec:s1]{Ingeniería de software}. \textbf{Fuentes:} \cite{scrum,swebok}. \hyperref[q:S-020]{Volver al reactivo}.\par\medskip

\subsection*{S-021 \quad Respuesta A}\label{s:S-021}
\textbf{Razonamiento.} La información visible reduce la dependencia de la memoria de trabajo.

\textbf{B) Distractor. }Aumenta la ambigüedad.

\textbf{C) Distractor. }Rompe consistencia y aprendizaje.

\textbf{Repaso:} \hyperref[sec:s1]{Ingeniería de software}. \textbf{Fuentes:} \cite{scrum,swebok}. \hyperref[q:S-021]{Volver al reactivo}.\par\medskip

\subsection*{S-022 \quad Respuesta B}\label{s:S-022}
\textbf{Razonamiento.} Se identifica el cambio, se valora y autoriza antes de ejecutarlo y comprobarlo.

\textbf{A) Distractor. }Aprueba sin registrar ni analizar la necesidad.

\textbf{C) Distractor. }Implementa antes de conocer alcance e impacto.

\textbf{Repaso:} \hyperref[sec:s1]{Ingeniería de software}. \textbf{Fuentes:} \cite{scrum,swebok}. \hyperref[q:S-022]{Volver al reactivo}.\par\medskip

\subsection*{S-023 \quad Respuesta C}\label{s:S-023}
\textbf{Razonamiento.} Cada nivel enfoca una perspectiva diferente de corrección.

\textbf{A) Distractor. }Asigna el contexto de usuario a la unidad y desplaza los demás.

\textbf{B) Distractor. }No representa los alcances habituales.

\textbf{Repaso:} \hyperref[sec:s1]{Ingeniería de software}. \textbf{Fuentes:} \cite{scrum,swebok}. \hyperref[q:S-023]{Volver al reactivo}.\par\medskip

\subsection*{S-024 \quad Respuesta B}\label{s:S-024}
\textbf{Razonamiento.} Las pruebas existentes pueden no representar la necesidad real o el flujo completo.

\textbf{A) Distractor. }Las pruebas siempre tienen alcance limitado.

\textbf{C) Distractor. }Compilar no detecta necesariamente el problema de uso.

\textbf{Repaso:} \hyperref[sec:s1]{Ingeniería de software}. \textbf{Fuentes:} \cite{scrum,swebok}. \hyperref[q:S-024]{Volver al reactivo}.\par\medskip

\subsection*{S-025 \quad Respuesta C}\label{s:S-025}
\textbf{Razonamiento.} Motivos de cambio independientes sugieren límites de responsabilidad más claros.

\textbf{A) Distractor. }Las responsabilidades no están enfocadas en una misma finalidad.

\textbf{B) Distractor. }Las dependencias deben analizarse, no presumirse inexistentes.

\textbf{Repaso:} \hyperref[sec:s1]{Ingeniería de software}. \textbf{Fuentes:} \cite{scrum,swebok}. \hyperref[q:S-025]{Volver al reactivo}.\par\medskip

\subsection*{S-026 \quad Respuesta A}\label{s:S-026}
\textbf{Razonamiento.} Un contrato consumido externamente requiere gestionar el impacto del cambio.

\textbf{B) Distractor. }La compilación interna no protege a los consumidores.

\textbf{C) Distractor. }La documentación es necesaria para adaptar las integraciones.

\textbf{Repaso:} \hyperref[sec:s1]{Ingeniería de software}. \textbf{Fuentes:} \cite{scrum,swebok}. \hyperref[q:S-026]{Volver al reactivo}.\par\medskip

\subsection*{S-027 \quad Respuesta C}\label{s:S-027}
\textbf{Razonamiento.} La estimación debe relacionarse con alcance, experiencia y riesgos.

\textbf{A) Distractor. }La precisión aparente no aporta sustento.

\textbf{B) Distractor. }El objetivo comercial no determina el esfuerzo real.

\textbf{Repaso:} \hyperref[sec:s1]{Ingeniería de software}. \textbf{Fuentes:} \cite{scrum,swebok}. \hyperref[q:S-027]{Volver al reactivo}.\par\medskip

\subsection*{S-028 \quad Respuesta A}\label{s:S-028}
\textbf{Razonamiento.} La cobertura indica ejecución, no suficiencia de comprobaciones ni de combinaciones.

\textbf{B) Distractor. }No es una demostración de ausencia de errores.

\textbf{C) Distractor. }Las combinaciones de entrada suelen superar los casos ejecutados.

\textbf{Repaso:} \hyperref[sec:s1]{Ingeniería de software}. \textbf{Fuentes:} \cite{scrum,swebok}. \hyperref[q:S-028]{Volver al reactivo}.\par\medskip

\subsection*{S-029 \quad Respuesta B}\label{s:S-029}
\textbf{Razonamiento.} Los criterios explícitos permiten justificar adecuación y consecuencias.

\textbf{A) Distractor. }La novedad no demuestra ajuste a las restricciones.

\textbf{C) Distractor. }Más clases no implica mejor resultado.

\textbf{Repaso:} \hyperref[sec:s1]{Ingeniería de software}. \textbf{Fuentes:} \cite{scrum,swebok}. \hyperref[q:S-029]{Volver al reactivo}.\par\medskip

\subsection*{S-030 \quad Respuesta A}\label{s:S-030}
\textbf{Razonamiento.} La incertidumbre relevante requiere seguimiento y una respuesta preparada.

\textbf{B) Distractor. }Impide actuar oportunamente.

\textbf{C) Distractor. }La ausencia de información no demuestra ausencia de riesgo.

\textbf{Repaso:} \hyperref[sec:s1]{Ingeniería de software}. \textbf{Fuentes:} \cite{scrum,swebok}. \hyperref[q:S-030]{Volver al reactivo}.\par\medskip

\subsection*{S-031 \quad Respuesta B}\label{s:S-031}
\textbf{Razonamiento.} Reducir datos innecesarios limita exposición y complejidad operativa.

\textbf{A) Distractor. }Amplía la exposición y dificulta el control.

\textbf{C) Distractor. }Aumenta el tratamiento sin justificación funcional.

\textbf{Repaso:} \hyperref[sec:s1]{Ingeniería de software}. \textbf{Fuentes:} \cite{scrum,swebok}. \hyperref[q:S-031]{Volver al reactivo}.\par\medskip

\subsection*{S-032 \quad Respuesta C}\label{s:S-032}
\textbf{Razonamiento.} La regresión protege contra la reintroducción del comportamiento defectuoso.

\textbf{A) Distractor. }No ejercita la causa ni el resultado del fallo.

\textbf{B) Distractor. }Escribir código no demuestra que el defecto no reaparezca.

\textbf{Repaso:} \hyperref[sec:s1]{Ingeniería de software}. \textbf{Fuentes:} \cite{scrum,swebok}. \hyperref[q:S-032]{Volver al reactivo}.\par\medskip

\section{Programación y software base}

\subsection*{S-033 \quad Respuesta C}\label{s:S-033}
\textbf{Razonamiento.} El consumidor depende de operaciones comunes y delega el comportamiento concreto.

\textbf{A) Distractor. }No requiere replicación integral.

\textbf{B) Distractor. }La sustitución compatible es precisamente una utilidad del contrato.

\textbf{Repaso:} \hyperref[sec:s2]{Programación y software base}. \textbf{Fuentes:} \cite{ostep,cppreference,compiler}. \hyperref[q:S-033]{Volver al reactivo}.\par\medskip

\subsection*{S-034 \quad Respuesta A}\label{s:S-034}
\textbf{Razonamiento.} El proceso incluye contexto de ejecución y recursos administrados.

\textbf{B) Distractor. }El archivo puede originar un proceso, pero no es su ejecución.

\textbf{C) Distractor. }No corresponde al concepto del sistema operativo.

\textbf{Repaso:} \hyperref[sec:s2]{Programación y software base}. \textbf{Fuentes:} \cite{ostep,cppreference,compiler}. \hyperref[q:S-034]{Volver al reactivo}.\par\medskip

\subsection*{S-035 \quad Respuesta C}\label{s:S-035}
\textbf{Razonamiento.} Los hilos comparten recursos del proceso y mantienen estado de ejecución propio.

\textbf{A) Distractor. }El sistema planifica su ejecución.

\textbf{B) Distractor. }No requieren almacenamiento físico exclusivo.

\textbf{Repaso:} \hyperref[sec:s2]{Programación y software base}. \textbf{Fuentes:} \cite{ostep,cppreference,compiler}. \hyperref[q:S-035]{Volver al reactivo}.\par\medskip

\subsection*{S-036 \quad Respuesta A}\label{s:S-036}
\textbf{Razonamiento.} El lexer reconoce unidades como identificadores, palabras clave y literales.

\textbf{B) Distractor. }La administración de memoria es otra tarea.

\textbf{C) Distractor. }La transmisión de red no es una fase de compilación.

\textbf{Repaso:} \hyperref[sec:s2]{Programación y software base}. \textbf{Fuentes:} \cite{ostep,cppreference,compiler}. \hyperref[q:S-036]{Volver al reactivo}.\par\medskip

\subsection*{S-037 \quad Respuesta B}\label{s:S-037}
\textbf{Razonamiento.} HTML describe estructura y semántica del contenido web.

\textbf{A) Distractor. }No asigna tiempo de procesador.

\textbf{C) Distractor. }No protege información mediante claves.

\textbf{Repaso:} \hyperref[sec:s2]{Programación y software base}. \textbf{Fuentes:} \cite{ostep,cppreference,compiler}. \hyperref[q:S-037]{Volver al reactivo}.\par\medskip

\subsection*{S-038 \quad Respuesta A}\label{s:S-038}
\textbf{Razonamiento.} El orden no controlado de operaciones compartidas puede alterar el resultado.

\textbf{B) Distractor. }El ordenamiento de filas no define concurrencia insegura.

\textbf{C) Distractor. }Puede haber concurrencia intercalada incluso en un solo núcleo.

\textbf{Repaso:} \hyperref[sec:s2]{Programación y software base}. \textbf{Fuentes:} \cite{ostep,cppreference,compiler}. \hyperref[q:S-038]{Volver al reactivo}.\par\medskip

\subsection*{S-039 \quad Respuesta B}\label{s:S-039}
\textbf{Razonamiento.} El sistema traduce direcciones y gestiona aislamiento y residencia de páginas.

\textbf{A) Distractor. }Los fallos de página son parte posible del mecanismo.

\textbf{C) Distractor. }El respaldo en disco es mucho más lento.

\textbf{Repaso:} \hyperref[sec:s2]{Programación y software base}. \textbf{Fuentes:} \cite{ostep,cppreference,compiler}. \hyperref[q:S-039]{Volver al reactivo}.\par\medskip

\subsection*{S-040 \quad Respuesta C}\label{s:S-040}
\textbf{Razonamiento.} Limita dependencias sobre detalles internos y preserva invariantes.

\textbf{A) Distractor. }Expone directamente la representación.

\textbf{B) Distractor. }La interfaz permite usar el componente.

\textbf{Repaso:} \hyperref[sec:s2]{Programación y software base}. \textbf{Fuentes:} \cite{ostep,cppreference,compiler}. \hyperref[q:S-040]{Volver al reactivo}.\par\medskip

\subsection*{S-041 \quad Respuesta B}\label{s:S-041}
\textbf{Razonamiento.} Primero x=5 y después x=5×4=20.

\textbf{A) Distractor. }No sigue las actualizaciones secuenciales indicadas.

\textbf{C) Distractor. }Trata las operaciones como una única expresión distinta.

\textbf{Repaso:} \hyperref[sec:s2]{Programación y software base}. \textbf{Fuentes:} \cite{ostep,cppreference,compiler}. \hyperref[q:S-041]{Volver al reactivo}.\par\medskip

\subsection*{S-042 \quad Respuesta C}\label{s:S-042}
\textbf{Razonamiento.} Se acumula 1+2+3+4=10.

\textbf{A) Distractor. }Corresponde a detenerse antes de incluir 4.

\textbf{B) Distractor. }Es el número de iteraciones, no la suma.

\textbf{Repaso:} \hyperref[sec:s2]{Programación y software base}. \textbf{Fuentes:} \cite{ostep,cppreference,compiler}. \hyperref[q:S-042]{Volver al reactivo}.\par\medskip

\subsection*{S-043 \quad Respuesta A}\label{s:S-043}
\textbf{Razonamiento.} B espera a que termine A; su propia ejecución dura después 5 ms.

\textbf{B) Distractor. }Es el tiempo de servicio de B.

\textbf{C) Distractor. }Es el tiempo de finalización de B desde t=0.

\textbf{Repaso:} \hyperref[sec:s2]{Programación y software base}. \textbf{Fuentes:} \cite{ostep,cppreference,compiler}. \hyperref[q:S-043]{Volver al reactivo}.\par\medskip

\subsection*{S-044 \quad Respuesta C}\label{s:S-044}
\textbf{Razonamiento.} A usa 0–2, B ejecuta 2–3 y termina; A continúa después.

\textbf{A) Distractor. }Ignora el primer turno de A.

\textbf{B) Distractor. }Es la finalización posterior de A.

\textbf{Repaso:} \hyperref[sec:s2]{Programación y software base}. \textbf{Fuentes:} \cite{ostep,cppreference,compiler}. \hyperref[q:S-044]{Volver al reactivo}.\par\medskip

\subsection*{S-045 \quad Respuesta A}\label{s:S-045}
\textbf{Razonamiento.} 4096=2\textasciicircum{}12 bytes; se necesitan 12 bits para seleccionar el byte.

\textbf{B) Distractor. }Confunde el valor en KiB con un número de bits.

\textbf{C) Distractor. }Es una longitud posible de dirección completa.

\textbf{Repaso:} \hyperref[sec:s2]{Programación y software base}. \textbf{Fuentes:} \cite{ostep,cppreference,compiler}. \hyperref[q:S-045]{Volver al reactivo}.\par\medskip

\subsection*{S-046 \quad Respuesta B}\label{s:S-046}
\textbf{Razonamiento.} La exclusión requiere que los participantes usen consistentemente la protección.

\textbf{A) Distractor. }volatile no proporciona exclusión mutua.

\textbf{C) Distractor. }El escritor podría modificar durante la lectura protegida.

\textbf{Repaso:} \hyperref[sec:s2]{Programación y software base}. \textbf{Fuentes:} \cite{ostep,cppreference,compiler}. \hyperref[q:S-046]{Volver al reactivo}.\par\medskip

\subsection*{S-047 \quad Respuesta A}\label{s:S-047}
\textbf{Razonamiento.} Un orden único elimina la espera circular de ese patrón.

\textbf{B) Distractor. }No elimina la condición que produce el bloqueo.

\textbf{C) Distractor. }Mantiene el deadlock.

\textbf{Repaso:} \hyperref[sec:s2]{Programación y software base}. \textbf{Fuentes:} \cite{ostep,cppreference,compiler}. \hyperref[q:S-047]{Volver al reactivo}.\par\medskip

\subsection*{S-048 \quad Respuesta B}\label{s:S-048}
\textbf{Razonamiento.} Un identificador estable permite reconocer el mismo intento lógico y no aplicarlo dos veces.

\textbf{A) Distractor. }Los reintentos son posibles en redes reales.

\textbf{C) Distractor. }La apariencia no controla el efecto del servidor.

\textbf{Repaso:} \hyperref[sec:s2]{Programación y software base}. \textbf{Fuentes:} \cite{ostep,cppreference,compiler}. \hyperref[q:S-048]{Volver al reactivo}.\par\medskip

\subsection*{S-049 \quad Respuesta C}\label{s:S-049}
\textbf{Razonamiento.} El cliente puede manipularse y no constituye una frontera confiable.

\textbf{A) Distractor. }Puede omitirse o alterarse la validación local.

\textbf{B) Distractor. }El texto informativo no impone reglas.

\textbf{Repaso:} \hyperref[sec:s2]{Programación y software base}. \textbf{Fuentes:} \cite{ostep,cppreference,compiler}. \hyperref[q:S-049]{Volver al reactivo}.\par\medskip

\subsection*{S-050 \quad Respuesta B}\label{s:S-050}
\textbf{Razonamiento.} El flujo léxico, sintáctico y semántico prepara una representación para generación.

\textbf{A) Distractor. }La comprobación semántica necesita una estructura ya analizada.

\textbf{C) Distractor. }Genera antes de conocer la estructura y significado.

\textbf{Repaso:} \hyperref[sec:s2]{Programación y software base}. \textbf{Fuentes:} \cite{ostep,cppreference,compiler}. \hyperref[q:S-050]{Volver al reactivo}.\par\medskip

\subsection*{S-051 \quad Respuesta C}\label{s:S-051}
\textbf{Razonamiento.} Los conceptos corresponden al modelo habitual de organización de un proceso.

\textbf{A) Distractor. }Mezcla instrucciones, invocaciones y memoria dinámica.

\textbf{B) Distractor. }No describe las funciones de las regiones.

\textbf{Repaso:} \hyperref[sec:s2]{Programación y software base}. \textbf{Fuentes:} \cite{ostep,cppreference,compiler}. \hyperref[q:S-051]{Volver al reactivo}.\par\medskip

\subsection*{S-052 \quad Respuesta A}\label{s:S-052}
\textbf{Razonamiento.} La presión de memoria provoca paginación excesiva y poco progreso útil.

\textbf{B) Distractor. }El patrón refleja insuficiencia de residencia o mala localidad.

\textbf{C) Distractor. }El intercambio genera mucha entrada/salida.

\textbf{Repaso:} \hyperref[sec:s2]{Programación y software base}. \textbf{Fuentes:} \cite{ostep,cppreference,compiler}. \hyperref[q:S-052]{Volver al reactivo}.\par\medskip

\subsection*{S-053 \quad Respuesta C}\label{s:S-053}
\textbf{Razonamiento.} La pila permite recordar profundidad no acotada en el modelo formal.

\textbf{A) Distractor. }Un lenguaje de balanceo arbitrario no es regular.

\textbf{B) Distractor. }Una cantidad fija de estados no distingue profundidad ilimitada.

\textbf{Repaso:} \hyperref[sec:s2]{Programación y software base}. \textbf{Fuentes:} \cite{ostep,cppreference,compiler}. \hyperref[q:S-053]{Volver al reactivo}.\par\medskip

\subsection*{S-054 \quad Respuesta A}\label{s:S-054}
\textbf{Razonamiento.} El acceso ocurre cuando la vida útil del objeto ya terminó.

\textbf{B) Distractor. }No se plantea un algoritmo recursivo.

\textbf{C) Distractor. }Puede afectar memoria y seguridad, no solo presentación.

\textbf{Repaso:} \hyperref[sec:s2]{Programación y software base}. \textbf{Fuentes:} \cite{ostep,cppreference,compiler}. \hyperref[q:S-054]{Volver al reactivo}.\par\medskip

\subsection*{S-055 \quad Respuesta B}\label{s:S-055}
\textbf{Razonamiento.} La espera de red debe permitir atender eventos de la interfaz.

\textbf{A) Distractor. }No libera el hilo que procesa interacción.

\textbf{C) Distractor. }Mantiene ocupado el hilo y empeora respuesta.

\textbf{Repaso:} \hyperref[sec:s2]{Programación y software base}. \textbf{Fuentes:} \cite{ostep,cppreference,compiler}. \hyperref[q:S-055]{Volver al reactivo}.\par\medskip

\subsection*{S-056 \quad Respuesta A}\label{s:S-056}
\textbf{Razonamiento.} La solicitud contextual explica necesidad y reduce privilegios innecesarios.

\textbf{B) Distractor. }Excede el uso descrito.

\textbf{C) Distractor. }Se debe manejar la negativa sin falsear datos.

\textbf{Repaso:} \hyperref[sec:s2]{Programación y software base}. \textbf{Fuentes:} \cite{ostep,cppreference,compiler}. \hyperref[q:S-056]{Volver al reactivo}.\par\medskip

\subsection*{S-057 \quad Respuesta B}\label{s:S-057}
\textbf{Razonamiento.} Se debe limitar la demanda pendiente y comunicar saturación.

\textbf{A) Distractor. }Los hilos también consumen recursos limitados.

\textbf{C) Distractor. }La memoria física es finita.

\textbf{Repaso:} \hyperref[sec:s2]{Programación y software base}. \textbf{Fuentes:} \cite{ostep,cppreference,compiler}. \hyperref[q:S-057]{Volver al reactivo}.\par\medskip

\subsection*{S-058 \quad Respuesta C}\label{s:S-058}
\textbf{Razonamiento.} Un timeout no determina el resultado remoto; hace falta resolver la incertidumbre.

\textbf{A) Distractor. }La respuesta puede haberse perdido después de confirmar el efecto.

\textbf{B) Distractor. }Puede duplicar cobros u otras acciones críticas.

\textbf{Repaso:} \hyperref[sec:s2]{Programación y software base}. \textbf{Fuentes:} \cite{ostep,cppreference,compiler}. \hyperref[q:S-058]{Volver al reactivo}.\par\medskip

\section{Manejo de datos y algoritmos}

\subsection*{S-059 \quad Respuesta B}\label{s:S-059}
\textbf{Razonamiento.} El último elemento insertado es el primero en salir.

\textbf{A) Distractor. }Un conjunto no define esa disciplina.

\textbf{C) Distractor. }FIFO extrae primero el más antiguo.

\textbf{Repaso:} \hyperref[sec:s3]{Manejo de datos y algoritmos}. \textbf{Fuentes:} \cite{algorithms,postgres,ai}. \hyperref[q:S-059]{Volver al reactivo}.\par\medskip

\subsection*{S-060 \quad Respuesta C}\label{s:S-060}
\textbf{Razonamiento.} La cola extrae en orden de llegada.

\textbf{A) Distractor. }La pila invierte el orden.

\textbf{B) Distractor. }El árbol ordena por claves y no por llegada.

\textbf{Repaso:} \hyperref[sec:s3]{Manejo de datos y algoritmos}. \textbf{Fuentes:} \cite{algorithms,postgres,ai}. \hyperref[q:S-060]{Volver al reactivo}.\par\medskip

\subsection*{S-061 \quad Respuesta B}\label{s:S-061}
\textbf{Razonamiento.} El orden permite descartar la mitad sin perder la clave.

\textbf{A) Distractor. }Funciona con longitudes pares o impares.

\textbf{C) Distractor. }Puede tener valores diferentes.

\textbf{Repaso:} \hyperref[sec:s3]{Manejo de datos y algoritmos}. \textbf{Fuentes:} \cite{algorithms,postgres,ai}. \hyperref[q:S-061]{Volver al reactivo}.\par\medskip

\subsection*{S-062 \quad Respuesta C}\label{s:S-062}
\textbf{Razonamiento.} La notación expresa una relación de crecimiento, no una medida física exacta.

\textbf{A) Distractor. }Depende de implementación y máquina, no de la notación sola.

\textbf{B) Distractor. }No define automáticamente el espacio usado.

\textbf{Repaso:} \hyperref[sec:s3]{Manejo de datos y algoritmos}. \textbf{Fuentes:} \cite{algorithms,postgres,ai}. \hyperref[q:S-062]{Volver al reactivo}.\par\medskip

\subsection*{S-063 \quad Respuesta A}\label{s:S-063}
\textbf{Razonamiento.} La clave impone unicidad e identificación de las filas.

\textbf{B) Distractor. }Viola las restricciones de una clave primaria.

\textbf{C) Distractor. }Su posición visual no define su función.

\textbf{Repaso:} \hyperref[sec:s3]{Manejo de datos y algoritmos}. \textbf{Fuentes:} \cite{algorithms,postgres,ai}. \hyperref[q:S-063]{Volver al reactivo}.\par\medskip

\subsection*{S-064 \quad Respuesta C}\label{s:S-064}
\textbf{Razonamiento.} La restricción relaciona valores y evita referencias inválidas según sus reglas.

\textbf{A) Distractor. }No exige duplicar toda la información.

\textbf{B) Distractor. }No determina el orden de presentación.

\textbf{Repaso:} \hyperref[sec:s3]{Manejo de datos y algoritmos}. \textbf{Fuentes:} \cite{algorithms,postgres,ai}. \hyperref[q:S-064]{Volver al reactivo}.\par\medskip

\subsection*{S-065 \quad Respuesta A}\label{s:S-065}
\textbf{Razonamiento.} Evita dejar una parte del cambio aplicado y otra no.

\textbf{B) Distractor. }Contradice la propiedad todo o nada.

\textbf{C) Distractor. }La atomicidad puede implementarse con almacenamiento persistente.

\textbf{Repaso:} \hyperref[sec:s3]{Manejo de datos y algoritmos}. \textbf{Fuentes:} \cite{algorithms,postgres,ai}. \hyperref[q:S-065]{Volver al reactivo}.\par\medskip

\subsection*{S-066 \quad Respuesta B}\label{s:S-066}
\textbf{Razonamiento.} La durabilidad mantiene los efectos confirmados ante fallas cubiertas por el diseño.

\textbf{A) Distractor. }El orden de texto no define persistencia.

\textbf{C) Distractor. }Aislamiento regula interacciones concurrentes.

\textbf{Repaso:} \hyperref[sec:s3]{Manejo de datos y algoritmos}. \textbf{Fuentes:} \cite{algorithms,postgres,ai}. \hyperref[q:S-066]{Volver al reactivo}.\par\medskip

\subsection*{S-067 \quad Respuesta A}\label{s:S-067}
\textbf{Razonamiento.} El modelo aprende una relación con salidas conocidas durante entrenamiento.

\textbf{B) Distractor. }El aprendizaje requiere información de entrenamiento.

\textbf{C) Distractor. }Puede medirse frente a etiquetas de validación.

\textbf{Repaso:} \hyperref[sec:s3]{Manejo de datos y algoritmos}. \textbf{Fuentes:} \cite{algorithms,postgres,ai}. \hyperref[q:S-067]{Volver al reactivo}.\par\medskip

\subsection*{S-068 \quad Respuesta B}\label{s:S-068}
\textbf{Razonamiento.} Orienta la exploración usando una estimación del problema.

\textbf{A) Distractor. }No corresponde a una función de búsqueda.

\textbf{C) Distractor. }Puede aproximar sin ser exacta.

\textbf{Repaso:} \hyperref[sec:s3]{Manejo de datos y algoritmos}. \textbf{Fuentes:} \cite{algorithms,postgres,ai}. \hyperref[q:S-068]{Volver al reactivo}.\par\medskip

\subsection*{S-069 \quad Respuesta C}\label{s:S-069}
\textbf{Razonamiento.} Permite detener el proceso al alcanzar una instancia resoluble directamente.

\textbf{A) Distractor. }La repetición sin avance puede no terminar.

\textbf{B) Distractor. }El tamaño del archivo no define la condición algorítmica.

\textbf{Repaso:} \hyperref[sec:s3]{Manejo de datos y algoritmos}. \textbf{Fuentes:} \cite{algorithms,postgres,ai}. \hyperref[q:S-069]{Volver al reactivo}.\par\medskip

\subsection*{S-070 \quad Respuesta B}\label{s:S-070}
\textbf{Razonamiento.} La normalización se apoya en claves y dependencias funcionales.

\textbf{A) Distractor. }Cambiar capitalización es una transformación textual distinta.

\textbf{C) Distractor. }Las relaciones son la base del modelo.

\textbf{Repaso:} \hyperref[sec:s3]{Manejo de datos y algoritmos}. \textbf{Fuentes:} \cite{algorithms,postgres,ai}. \hyperref[q:S-070]{Volver al reactivo}.\par\medskip

\subsection*{S-071 \quad Respuesta C}\label{s:S-071}
\textbf{Razonamiento.} Se quita 5, se agrega 7 y este queda en la cima.

\textbf{A) Distractor. }5 ya se extrajo en el primer pop.

\textbf{B) Distractor. }2 permanece debajo de 7.

\textbf{Repaso:} \hyperref[sec:s3]{Manejo de datos y algoritmos}. \textbf{Fuentes:} \cite{algorithms,postgres,ai}. \hyperref[q:S-071]{Volver al reactivo}.\par\medskip

\subsection*{S-072 \quad Respuesta A}\label{s:S-072}
\textbf{Razonamiento.} A sale primero; quedan B y luego C.

\textbf{B) Distractor. }C llegó después de B.

\textbf{C) Distractor. }A ya fue retirado.

\textbf{Repaso:} \hyperref[sec:s3]{Manejo de datos y algoritmos}. \textbf{Fuentes:} \cite{algorithms,postgres,ai}. \hyperref[q:S-072]{Volver al reactivo}.\par\medskip

\subsection*{S-073 \quad Respuesta C}\label{s:S-073}
\textbf{Razonamiento.} 1024=2\textasciicircum{}10; diez divisiones entre dos dejan tamaño uno.

\textbf{A) Distractor. }Describe un recorrido lineal.

\textbf{B) Distractor. }No coincide con el logaritmo en base dos.

\textbf{Repaso:} \hyperref[sec:s3]{Manejo de datos y algoritmos}. \textbf{Fuentes:} \cite{algorithms,postgres,ai}. \hyperref[q:S-073]{Volver al reactivo}.\par\medskip

\subsection*{S-074 \quad Respuesta A}\label{s:S-074}
\textbf{Razonamiento.} El cuerpo se ejecuta n×n veces.

\textbf{B) Distractor. }Solo cuenta uno de los bucles.

\textbf{C) Distractor. }No hay reducción geométrica del problema.

\textbf{Repaso:} \hyperref[sec:s3]{Manejo de datos y algoritmos}. \textbf{Fuentes:} \cite{algorithms,postgres,ai}. \hyperref[q:S-074]{Volver al reactivo}.\par\medskip

\subsection*{S-075 \quad Respuesta B}\label{s:S-075}
\textbf{Razonamiento.} Debe seguir los enlaces desde la cabeza hasta el final.

\textbf{A) Distractor. }No puede descartar mitades como en búsqueda binaria.

\textbf{C) Distractor. }No hay acceso directo al último nodo.

\textbf{Repaso:} \hyperref[sec:s3]{Manejo de datos y algoritmos}. \textbf{Fuentes:} \cite{algorithms,postgres,ai}. \hyperref[q:S-075]{Volver al reactivo}.\par\medskip

\subsection*{S-076 \quad Respuesta A}\label{s:S-076}
\textbf{Razonamiento.} BFS explora por niveles de distancia desde el origen.

\textbf{B) Distractor. }DFS puede profundizar por un camino largo antes de otro corto.

\textbf{C) Distractor. }Es un algoritmo de ordenamiento de secuencias.

\textbf{Repaso:} \hyperref[sec:s3]{Manejo de datos y algoritmos}. \textbf{Fuentes:} \cite{algorithms,postgres,ai}. \hyperref[q:S-076]{Volver al reactivo}.\par\medskip

\subsection*{S-077 \quad Respuesta B}\label{s:S-077}
\textbf{Razonamiento.} La relajación con extracción de menor distancia es correcta bajo pesos no negativos.

\textbf{A) Distractor. }Los nombres de nodos no determinan los costos.

\textbf{C) Distractor. }No describe la función del algoritmo.

\textbf{Repaso:} \hyperref[sec:s3]{Manejo de datos y algoritmos}. \textbf{Fuentes:} \cite{algorithms,postgres,ai}. \hyperref[q:S-077]{Volver al reactivo}.\par\medskip

\subsection*{S-078 \quad Respuesta C}\label{s:S-078}
\textbf{Razonamiento.} Bajo esos supuestos el número esperado de inspecciones está acotado.

\textbf{A) Distractor. }No es el comportamiento esperado de una implementación adecuada.

\textbf{B) Distractor. }Hay que calcular hash y verificar coincidencias.

\textbf{Repaso:} \hyperref[sec:s3]{Manejo de datos y algoritmos}. \textbf{Fuentes:} \cite{algorithms,postgres,ai}. \hyperref[q:S-078]{Volver al reactivo}.\par\medskip

\subsection*{S-079 \quad Respuesta B}\label{s:S-079}
\textbf{Razonamiento.} 4×3×2×1=24.

\textbf{A) Distractor. }Es 4 al cuadrado, no factorial.

\textbf{C) Distractor. }Es la suma de 1 a 4.

\textbf{Repaso:} \hyperref[sec:s3]{Manejo de datos y algoritmos}. \textbf{Fuentes:} \cite{algorithms,postgres,ai}. \hyperref[q:S-079]{Volver al reactivo}.\par\medskip

\subsection*{S-080 \quad Respuesta C}\label{s:S-080}
\textbf{Razonamiento.} SUM agrega los valores no nulos: 10+20=30.

\textbf{A) Distractor. }La presencia de un NULL no anula todos los sumandos.

\textbf{B) Distractor. }3 es el número total de filas, no la suma.

\textbf{Repaso:} \hyperref[sec:s3]{Manejo de datos y algoritmos}. \textbf{Fuentes:} \cite{postgres}. \hyperref[q:S-080]{Volver al reactivo}.\par\medskip

\subsection*{S-081 \quad Respuesta A}\label{s:S-081}
\textbf{Razonamiento.} IS NULL prueba ausencia de valor.

\textbf{B) Distractor. }La comparación ordinaria con NULL produce desconocido.

\textbf{C) Distractor. }Compara con una cadena literal de cuatro letras.

\textbf{Repaso:} \hyperref[sec:s3]{Manejo de datos y algoritmos}. \textbf{Fuentes:} \cite{postgres}. \hyperref[q:S-081]{Volver al reactivo}.\par\medskip

\subsection*{S-082 \quad Respuesta C}\label{s:S-082}
\textbf{Razonamiento.} Solo se emiten las tres coincidencias existentes.

\textbf{A) Distractor. }No produce una sola fila por cliente si hay varias coincidencias.

\textbf{B) Distractor. }Seis sería el producto cartesiano sin condición adecuada.

\textbf{Repaso:} \hyperref[sec:s3]{Manejo de datos y algoritmos}. \textbf{Fuentes:} \cite{postgres}. \hyperref[q:S-082]{Volver al reactivo}.\par\medskip

\subsection*{S-083 \quad Respuesta A}\label{s:S-083}
\textbf{Razonamiento.} HAVING filtra resultados de agregación después de agrupar por cliente.

\textbf{B) Distractor. }WHERE no evalúa agregados de grupo en esa posición.

\textbf{C) Distractor. }Agrupa por un atributo diferente y no aplica el filtro solicitado.

\textbf{Repaso:} \hyperref[sec:s3]{Manejo de datos y algoritmos}. \textbf{Fuentes:} \cite{postgres}. \hyperref[q:S-083]{Volver al reactivo}.\par\medskip

\subsection*{S-084 \quad Respuesta B}\label{s:S-084}
\textbf{Razonamiento.} Un atributo no primo depende solo de una parte propia de la clave compuesta.

\textbf{A) Distractor. }La dependencia está expresamente dada.

\textbf{C) Distractor. }La dependencia a alumno, que no es superclave, impide afirmar BCNF.

\textbf{Repaso:} \hyperref[sec:s3]{Manejo de datos y algoritmos}. \textbf{Fuentes:} \cite{algorithms,postgres,ai}. \hyperref[q:S-084]{Volver al reactivo}.\par\medskip

\subsection*{S-085 \quad Respuesta A}\label{s:S-085}
\textbf{Razonamiento.} Debe mantenerse la estructura adicional cuando cambian los datos.

\textbf{B) Distractor. }El mantenimiento puede aumentar el trabajo de escritura.

\textbf{C) Distractor. }El optimizador elige según selectividad y costos.

\textbf{Repaso:} \hyperref[sec:s3]{Manejo de datos y algoritmos}. \textbf{Fuentes:} \cite{postgres}. \hyperref[q:S-085]{Volver al reactivo}.\par\medskip

\subsection*{S-086 \quad Respuesta B}\label{s:S-086}
\textbf{Razonamiento.} El resultado debería ser 80, pero una escritura reemplaza el efecto de la otra.

\textbf{A) Distractor. }La normalización no describe intercalación de escrituras.

\textbf{C) Distractor. }El saldo final de 90 omite un descuento.

\textbf{Repaso:} \hyperref[sec:s3]{Manejo de datos y algoritmos}. \textbf{Fuentes:} \cite{postgres}. \hyperref[q:S-086]{Volver al reactivo}.\par\medskip

\subsection*{S-087 \quad Respuesta C}\label{s:S-087}
\textbf{Razonamiento.} La atomicidad conjunta evita un débito aislado sin el crédito correspondiente.

\textbf{A) Distractor. }Puede confirmar solo la primera mitad.

\textbf{B) Distractor. }Leer no efectúa ni protege los dos cambios.

\textbf{Repaso:} \hyperref[sec:s3]{Manejo de datos y algoritmos}. \textbf{Fuentes:} \cite{postgres}. \hyperref[q:S-087]{Volver al reactivo}.\par\medskip

\subsection*{S-088 \quad Respuesta B}\label{s:S-088}
\textbf{Razonamiento.} Exactitud=predicciones correctas/total=90/100.

\textbf{A) Distractor. }No todas las predicciones fueron correctas.

\textbf{C) Distractor. }Es la tasa de error en ese conjunto.

\textbf{Repaso:} \hyperref[sec:s3]{Manejo de datos y algoritmos}. \textbf{Fuentes:} \cite{algorithms,postgres,ai}. \hyperref[q:S-088]{Volver al reactivo}.\par\medskip

\subsection*{S-089 \quad Respuesta C}\label{s:S-089}
\textbf{Razonamiento.} Recall=TP/(TP+FN)=15/20=0.75.

\textbf{A) Distractor. }Divide por un total no proporcionado de 100.

\textbf{B) Distractor. }Es la fracción de positivos omitidos.

\textbf{Repaso:} \hyperref[sec:s3]{Manejo de datos y algoritmos}. \textbf{Fuentes:} \cite{algorithms,postgres,ai}. \hyperref[q:S-089]{Volver al reactivo}.\par\medskip

\subsection*{S-090 \quad Respuesta A}\label{s:S-090}
\textbf{Razonamiento.} La brecha exige revisar complejidad, datos y separación de evaluación.

\textbf{B) Distractor. }La validación mala contradice la afirmación.

\textbf{C) Distractor. }No se deduce que falten todas las variables.

\textbf{Repaso:} \hyperref[sec:s3]{Manejo de datos y algoritmos}. \textbf{Fuentes:} \cite{algorithms,postgres,ai}. \hyperref[q:S-090]{Volver al reactivo}.\par\medskip

\subsection*{S-091 \quad Respuesta C}\label{s:S-091}
\textbf{Razonamiento.} La reserva se define antes de aprender y el conjunto de prueba no guía los ajustes.

\textbf{A) Distractor. }Separar después permite contaminar la evaluación.

\textbf{B) Distractor. }No se evalúa un modelo definitivo antes de entrenarlo.

\textbf{Repaso:} \hyperref[sec:s3]{Manejo de datos y algoritmos}. \textbf{Fuentes:} \cite{algorithms,postgres,ai}. \hyperref[q:S-091]{Volver al reactivo}.\par\medskip

\subsection*{S-092 \quad Respuesta A}\label{s:S-092}
\textbf{Razonamiento.} Son las disciplinas y estructura indicadas por sus definiciones.

\textbf{B) Distractor. }Intercambia extracción y organización por clave.

\textbf{C) Distractor. }No coincide con el comportamiento de las estructuras.

\textbf{Repaso:} \hyperref[sec:s3]{Manejo de datos y algoritmos}. \textbf{Fuentes:} \cite{algorithms,postgres,ai}. \hyperref[q:S-092]{Volver al reactivo}.\par\medskip

\subsection*{S-093 \quad Respuesta B}\label{s:S-093}
\textbf{Razonamiento.} Bellman-Ford admite pesos negativos y permite detectar ciclos negativos.

\textbf{A) Distractor. }Minimiza aristas, no necesariamente suma de pesos.

\textbf{C) Distractor. }La garantía usual de Dijkstra exige pesos no negativos.

\textbf{Repaso:} \hyperref[sec:s3]{Manejo de datos y algoritmos}. \textbf{Fuentes:} \cite{algorithms,postgres,ai}. \hyperref[q:S-093]{Volver al reactivo}.\par\medskip

\subsection*{S-094 \quad Respuesta A}\label{s:S-094}
\textbf{Razonamiento.} Su altura puede crecer hasta n, y la búsqueda recorre esa cadena.

\textbf{B) Distractor. }Requiere altura logarítmica, que aquí no está garantizada.

\textbf{C) Distractor. }No existe acceso directo garantizado a cualquier clave.

\textbf{Repaso:} \hyperref[sec:s3]{Manejo de datos y algoritmos}. \textbf{Fuentes:} \cite{algorithms,postgres,ai}. \hyperref[q:S-094]{Volver al reactivo}.\par\medskip

\subsection*{S-095 \quad Respuesta B}\label{s:S-095}
\textbf{Razonamiento.} La admisibilidad exige h(n)≤costo óptimo restante.

\textbf{A) Distractor. }No depende de ser aleatoria.

\textbf{C) Distractor. }Puede subestimar estrictamente.

\textbf{Repaso:} \hyperref[sec:s3]{Manejo de datos y algoritmos}. \textbf{Fuentes:} \cite{ai}. \hyperref[q:S-095]{Volver al reactivo}.\par\medskip

\subsection*{S-096 \quad Respuesta C}\label{s:S-096}
\textbf{Razonamiento.} El aislamiento puede rechazar intercalaciones; la unidad de reintento es la transacción lógica.

\textbf{A) Distractor. }El aborto significa que el cambio no se confirmó.

\textbf{B) Distractor. }Los resultados pueden depender de una instantánea que ya no permite la operación.

\textbf{Repaso:} \hyperref[sec:s3]{Manejo de datos y algoritmos}. \textbf{Fuentes:} \cite{postgres}. \hyperref[q:S-096]{Volver al reactivo}.\par\medskip

\subsection*{S-097 \quad Respuesta B}\label{s:S-097}
\textbf{Razonamiento.} Los valores de cabecera se repiten por cada coincidencia; la agregación debe respetar esa granularidad.

\textbf{A) Distractor. }Más memoria no corrige una suma semánticamente incorrecta.

\textbf{C) Distractor. }La presentación no causa la multiplicación lógica.

\textbf{Repaso:} \hyperref[sec:s3]{Manejo de datos y algoritmos}. \textbf{Fuentes:} \cite{postgres}. \hyperref[q:S-097]{Volver al reactivo}.\par\medskip

\subsection*{S-098 \quad Respuesta C}\label{s:S-098}
\textbf{Razonamiento.} La similitud compartida puede permitir memorizar en vez de generalizar.

\textbf{A) Distractor. }Los conjuntos dejan de aportar evidencia verdaderamente separada.

\textbf{B) Distractor. }La validez en nuevos casos sigue necesitando evaluación.

\textbf{Repaso:} \hyperref[sec:s3]{Manejo de datos y algoritmos}. \textbf{Fuentes:} \cite{algorithms,postgres,ai}. \hyperref[q:S-098]{Volver al reactivo}.\par\medskip

\subsection*{S-099 \quad Respuesta A}\label{s:S-099}
\textbf{Razonamiento.} La memoización o tabulación evita recomputar y usa la estructura óptima.

\textbf{B) Distractor. }Desperdicia la repetición disponible.

\textbf{C) Distractor. }No usa ninguna de las propiedades dadas.

\textbf{Repaso:} \hyperref[sec:s3]{Manejo de datos y algoritmos}. \textbf{Fuentes:} \cite{algorithms,postgres,ai}. \hyperref[q:S-099]{Volver al reactivo}.\par\medskip

\subsection*{S-100 \quad Respuesta C}\label{s:S-100}
\textbf{Razonamiento.} Un árbol balanceado limita la altura a O(log n), permite localizar el inicio y añade O(k) para devolver k coincidencias; la consulta completa no es solo logarítmica.

\textbf{A) Distractor. }Una tabla hash no mantiene por sí sola el orden de las claves.

\textbf{B) Distractor. }La búsqueda en lista no ordenada puede recorrer todos los elementos.

\textbf{Repaso:} \hyperref[sec:s3]{Manejo de datos y algoritmos}. \textbf{Fuentes:} \cite{algorithms,postgres,ai}. \hyperref[q:S-100]{Volver al reactivo}.\par\medskip

\section{Comprensión lectora: ámbito de estudio}

\subsection*{L-001 \quad Respuesta C}\label{s:L-001}
\textbf{Razonamiento.} El primer párrafo delimita explícitamente los talleres a dos grupos de secundaria.

\textbf{A) Distractor. }El texto no describe una muestra nacional.

\textbf{B) Distractor. }El puente ilustra un uso técnico, no identifica participantes.

\textbf{Repaso:} \hyperref[sec:l1]{Comprensión lectora: ámbito de estudio}. \textbf{Fuentes:} \cite{ceneval,reading}. \hyperref[q:L-001]{Volver al reactivo}.\par\medskip

\subsection*{L-002 \quad Respuesta A}\label{s:L-002}
\textbf{Razonamiento.} El primer párrafo enumera biblioteca, paso durante lluvias y punto de reunión.

\textbf{B) Distractor. }No se dice que se agregaran nombres de alumnos.

\textbf{C) Distractor. }El puente se menciona como contraste de propósito.

\textbf{Repaso:} \hyperref[sec:l1]{Comprensión lectora: ámbito de estudio}. \textbf{Fuentes:} \cite{ceneval,reading}. \hyperref[q:L-002]{Volver al reactivo}.\par\medskip

\subsection*{L-003 \quad Respuesta C}\label{s:L-003}
\textbf{Razonamiento.} El segundo párrafo identifica expresamente esa falta de detalle.

\textbf{A) Distractor. }La reseña afirma que la explicación de talleres es clara.

\textbf{B) Distractor. }La incorporación de esos recorridos es la propuesta del libro.

\textbf{Repaso:} \hyperref[sec:l1]{Comprensión lectora: ámbito de estudio}. \textbf{Fuentes:} \cite{ceneval,reading}. \hyperref[q:L-003]{Volver al reactivo}.\par\medskip

\subsection*{L-004 \quad Respuesta A}\label{s:L-004}
\textbf{Razonamiento.} La reseña reconoce un aporte y limita el alcance de las conclusiones.

\textbf{B) Distractor. }El texto rechaza explícitamente la sustitución para toda tarea.

\textbf{C) Distractor. }La limitación no lleva a negar toda utilidad.

\textbf{Repaso:} \hyperref[sec:l1]{Comprensión lectora: ámbito de estudio}. \textbf{Fuentes:} \cite{ceneval,reading}. \hyperref[q:L-004]{Volver al reactivo}.\par\medskip

\subsection*{L-005 \quad Respuesta B}\label{s:L-005}
\textbf{Razonamiento.} El ejemplo contrasta itinerarios personales y planos técnicos para tareas diferentes.

\textbf{A) Distractor. }La tesis central sostiene que sí hay selección según finalidad.

\textbf{C) Distractor. }No se relata ninguna construcción.

\textbf{Repaso:} \hyperref[sec:l1]{Comprensión lectora: ámbito de estudio}. \textbf{Fuentes:} \cite{ceneval,reading}. \hyperref[q:L-005]{Volver al reactivo}.\par\medskip

\subsection*{L-006 \quad Respuesta A}\label{s:L-006}
\textbf{Razonamiento.} La reseña dice que no basta para asegurar beneficios para todos los estudiantes.

\textbf{B) Distractor. }El primer párrafo reporta precisamente esa observación.

\textbf{C) Distractor. }Es una limitación que el texto sí formula.

\textbf{Repaso:} \hyperref[sec:l1]{Comprensión lectora: ámbito de estudio}. \textbf{Fuentes:} \cite{ceneval,reading}. \hyperref[q:L-006]{Volver al reactivo}.\par\medskip

\subsection*{L-007 \quad Respuesta B}\label{s:L-007}
\textbf{Razonamiento.} De las cuarenta personas, veinte fueron asignadas al grupo con recordatorio.

\textbf{A) Distractor. }Es el número de devoluciones puntuales en ese grupo.

\textbf{C) Distractor. }Es la muestra total.

\textbf{Repaso:} \hyperref[sec:l1]{Comprensión lectora: ámbito de estudio}. \textbf{Fuentes:} \cite{ceneval,reading}. \hyperref[q:L-007]{Volver al reactivo}.\par\medskip

\subsection*{L-008 \quad Respuesta C}\label{s:L-008}
\textbf{Razonamiento.} El primer párrafo aclara que ambos grupos conservaron las mismas condiciones.

\textbf{A) Distractor. }Los resultados fueron 16 y 12, no todos.

\textbf{B) Distractor. }El informe señala que no se registró apertura del aviso.

\textbf{Repaso:} \hyperref[sec:l1]{Comprensión lectora: ámbito de estudio}. \textbf{Fuentes:} \cite{ceneval,reading}. \hyperref[q:L-008]{Volver al reactivo}.\par\medskip

\subsection*{L-009 \quad Respuesta B}\label{s:L-009}
\textbf{Razonamiento.} 16/20=80\% y 12/20=60\%; la diferencia es 20 puntos.

\textbf{A) Distractor. }No coincide con las proporciones de los grupos.

\textbf{C) Distractor. }Cuatro es la diferencia en personas, no en puntos porcentuales.

\textbf{Repaso:} \hyperref[sec:l1]{Comprensión lectora: ámbito de estudio}. \textbf{Fuentes:} \cite{ceneval,reading}. \hyperref[q:L-009]{Volver al reactivo}.\par\medskip

\subsection*{L-010 \quad Respuesta C}\label{s:L-010}
\textbf{Razonamiento.} El segundo párrafo distingue resultado de muestra y generalización.

\textbf{A) Distractor. }El informe evita esa conclusión universal.

\textbf{B) Distractor. }La falta de registro de apertura no prueba que nadie lo abriera.

\textbf{Repaso:} \hyperref[sec:l1]{Comprensión lectora: ámbito de estudio}. \textbf{Fuentes:} \cite{ceneval,reading}. \hyperref[q:L-010]{Volver al reactivo}.\par\medskip

\subsection*{L-011 \quad Respuesta A}\label{s:L-011}
\textbf{Razonamiento.} El sorteo asigna la intervención sin elección intencional de participantes para cada grupo.

\textbf{B) Distractor. }No fuerza igualdad de resultados.

\textbf{C) Distractor. }El texto reconoce limitaciones pese al sorteo.

\textbf{Repaso:} \hyperref[sec:l1]{Comprensión lectora: ámbito de estudio}. \textbf{Fuentes:} \cite{ceneval,reading}. \hyperref[q:L-011]{Volver al reactivo}.\par\medskip

\subsection*{L-012 \quad Respuesta C}\label{s:L-012}
\textbf{Razonamiento.} El informe reconoce no haber registrado la apertura y propone medir recepción efectiva.

\textbf{A) Distractor. }Cambiar condiciones después altera la interpretación.

\textbf{B) Distractor. }Excluir casos relevantes sesga el resultado.

\textbf{Repaso:} \hyperref[sec:l1]{Comprensión lectora: ámbito de estudio}. \textbf{Fuentes:} \cite{ceneval,reading}. \hyperref[q:L-012]{Volver al reactivo}.\par\medskip

\section{Comprensión lectora: ámbito literario}

\subsection*{L-013 \quad Respuesta C}\label{s:L-013}
\textbf{Razonamiento.} El primer párrafo describe colocación y limpieza, y aclara que no se sentaba.

\textbf{A) Distractor. }No se relata que se sentara ni que el hermano llegara.

\textbf{B) Distractor. }No aparece ninguna venta.

\textbf{Repaso:} \hyperref[sec:l2]{Comprensión lectora: ámbito literario}. \textbf{Fuentes:} \cite{ceneval,reading}. \hyperref[q:L-013]{Volver al reactivo}.\par\medskip

\subsection*{L-014 \quad Respuesta A}\label{s:L-014}
\textbf{Razonamiento.} El segundo párrafo dice que reconoció la letra de su hermano.

\textbf{B) Distractor. }El niño aparece después y no es identificado como remitente.

\textbf{C) Distractor. }No se identifica a un vecino como autor.

\textbf{Repaso:} \hyperref[sec:l2]{Comprensión lectora: ámbito literario}. \textbf{Fuentes:} \cite{ceneval,reading}. \hyperref[q:L-014]{Volver al reactivo}.\par\medskip

\subsection*{L-015 \quad Respuesta C}\label{s:L-015}
\textbf{Razonamiento.} El segundo párrafo relata que la sacó de todos modos y se quedó a su lado.

\textbf{A) Distractor. }No se menciona que destruyera la silla.

\textbf{B) Distractor. }Contradice la salida a la lluvia.

\textbf{Repaso:} \hyperref[sec:l2]{Comprensión lectora: ámbito literario}. \textbf{Fuentes:} \cite{ceneval,reading}. \hyperref[q:L-015]{Volver al reactivo}.\par\medskip

\subsection*{L-016 \quad Respuesta A}\label{s:L-016}
\textbf{Razonamiento.} La aclaración de que no estaba rota desplaza el sentido del arreglo hacia lo simbólico.

\textbf{B) Distractor. }El narrador niega esa necesidad material.

\textbf{C) Distractor. }El niño no envía instrucciones.

\textbf{Repaso:} \hyperref[sec:l2]{Comprensión lectora: ámbito literario}. \textbf{Fuentes:} \cite{ceneval,reading}. \hyperref[q:L-016]{Volver al reactivo}.\par\medskip

\subsection*{L-017 \quad Respuesta B}\label{s:L-017}
\textbf{Razonamiento.} La frase y la apertura del taller indican una disponibilidad nueva.

\textbf{A) Distractor. }La invitación contradice la prohibición.

\textbf{C) Distractor. }El cuento no narra la llegada del hermano.

\textbf{Repaso:} \hyperref[sec:l2]{Comprensión lectora: ámbito literario}. \textbf{Fuentes:} \cite{ceneval,reading}. \hyperref[q:L-017]{Volver al reactivo}.\par\medskip

\subsection*{L-018 \quad Respuesta A}\label{s:L-018}
\textbf{Razonamiento.} La narración omite el texto y muestra cambios de conducta posteriores.

\textbf{B) Distractor. }El contenido no se transcribe.

\textbf{C) Distractor. }La omisión no demuestra ausencia de contenido.

\textbf{Repaso:} \hyperref[sec:l2]{Comprensión lectora: ámbito literario}. \textbf{Fuentes:} \cite{ceneval,reading}. \hyperref[q:L-018]{Volver al reactivo}.\par\medskip

\subsection*{L-019 \quad Respuesta B}\label{s:L-019}
\textbf{Razonamiento.} El primer párrafo vincula la mancha con la mesa de una estación.

\textbf{A) Distractor. }El narrador dice que no conserva una sola fecha.

\textbf{C) Distractor. }No se mencionan fotografías.

\textbf{Repaso:} \hyperref[sec:l2]{Comprensión lectora: ámbito literario}. \textbf{Fuentes:} \cite{ceneval,reading}. \hyperref[q:L-019]{Volver al reactivo}.\par\medskip

\subsection*{L-020 \quad Respuesta C}\label{s:L-020}
\textbf{Razonamiento.} El segundo párrafo introduce esa petición de cuidado.

\textbf{A) Distractor. }Es justamente el exceso que el ensayo problematiza.

\textbf{B) Distractor. }No es la solicitud atribuida a quien presta.

\textbf{Repaso:} \hyperref[sec:l2]{Comprensión lectora: ámbito literario}. \textbf{Fuentes:} \cite{ceneval,reading}. \hyperref[q:L-020]{Volver al reactivo}.\par\medskip

\subsection*{L-021 \quad Respuesta B}\label{s:L-021}
\textbf{Razonamiento.} El segundo párrafo sitúa el problema en que la voz propia ocupe demasiado espacio.

\textbf{A) Distractor. }El tercer párrafo rechaza convertir cada libro en un cuaderno de ejercicios.

\textbf{C) Distractor. }El ensayo distingue marcas de daños.

\textbf{Repaso:} \hyperref[sec:l2]{Comprensión lectora: ámbito literario}. \textbf{Fuentes:} \cite{ceneval,reading}. \hyperref[q:L-021]{Volver al reactivo}.\par\medskip

\subsection*{L-022 \quad Respuesta C}\label{s:L-022}
\textbf{Razonamiento.} El tercer párrafo explica la orilla como distancia habitable entre lector y texto.

\textbf{A) Distractor. }La imagen es reflexiva, no una instrucción de almacenamiento.

\textbf{B) Distractor. }El ensayo sí admite anotaciones como conversación.

\textbf{Repaso:} \hyperref[sec:l2]{Comprensión lectora: ámbito literario}. \textbf{Fuentes:} \cite{ceneval,reading}. \hyperref[q:L-022]{Volver al reactivo}.\par\medskip

\subsection*{L-023 \quad Respuesta A}\label{s:L-023}
\textbf{Razonamiento.} El autor defiende una práctica moderada y dependiente de la ocasión.

\textbf{B) Distractor. }Reconoce el valor de algunas anotaciones.

\textbf{C) Distractor. }También considera valioso cerrar el lápiz y seguir leyendo.

\textbf{Repaso:} \hyperref[sec:l2]{Comprensión lectora: ámbito literario}. \textbf{Fuentes:} \cite{ceneval,reading}. \hyperref[q:L-023]{Volver al reactivo}.\par\medskip

\subsection*{L-024 \quad Respuesta C}\label{s:L-024}
\textbf{Razonamiento.} La anécdota abre el tema que después se desarrolla en forma reflexiva.

\textbf{A) Distractor. }No hay muestra ni inferencia estadística.

\textbf{B) Distractor. }No se formula esa instrucción normativa.

\textbf{Repaso:} \hyperref[sec:l2]{Comprensión lectora: ámbito literario}. \textbf{Fuentes:} \cite{ceneval,reading}. \hyperref[q:L-024]{Volver al reactivo}.\par\medskip

\section{Comprensión lectora: participación social}

\subsection*{L-025 \quad Respuesta C}\label{s:L-025}
\textbf{Razonamiento.} Cumple residencia, edad, extensión y entrega dentro del periodo indicado.

\textbf{A) Distractor. }No supera los dieciséis años exigidos.

\textbf{B) Distractor. }La extensión supera el máximo de 500 palabras.

\textbf{Repaso:} \hyperref[sec:l3]{Comprensión lectora: participación social}. \textbf{Fuentes:} \cite{ceneval,reading}. \hyperref[q:L-025]{Volver al reactivo}.\par\medskip

\subsection*{L-026 \quad Respuesta A}\label{s:L-026}
\textbf{Razonamiento.} El segundo párrafo define sorteo cuando la demanda válida supera los lugares.

\textbf{B) Distractor. }La convocatoria indica que no se requiere experiencia.

\textbf{C) Distractor. }No establece una regla alfabética.

\textbf{Repaso:} \hyperref[sec:l3]{Comprensión lectora: participación social}. \textbf{Fuentes:} \cite{ceneval,reading}. \hyperref[q:L-026]{Volver al reactivo}.\par\medskip

\subsection*{L-027 \quad Respuesta C}\label{s:L-027}
\textbf{Razonamiento.} El cierre de la convocatoria vincula el aviso con liberar el lugar.

\textbf{A) Distractor. }No se establece sanción económica.

\textbf{B) Distractor. }No se prevé cancelar el taller por esa circunstancia.

\textbf{Repaso:} \hyperref[sec:l3]{Comprensión lectora: participación social}. \textbf{Fuentes:} \cite{ceneval,reading}. \hyperref[q:L-027]{Volver al reactivo}.\par\medskip

\subsection*{L-028 \quad Respuesta A}\label{s:L-028}
\textbf{Razonamiento.} El primer párrafo delimita la intervención a cuatro calles.

\textbf{B) Distractor. }Aclara que no cambia el servicio en otras colonias.

\textbf{C) Distractor. }El mercado aparece por su asociación, no como único ámbito.

\textbf{Repaso:} \hyperref[sec:l3]{Comprensión lectora: participación social}. \textbf{Fuentes:} \cite{ceneval,reading}. \hyperref[q:L-028]{Volver al reactivo}.\par\medskip

\subsection*{L-029 \quad Respuesta B}\label{s:L-029}
\textbf{Razonamiento.} El segundo párrafo afirma que aún no se ha decidido ampliar.

\textbf{A) Distractor. }No se plantea inscripción ni ese requisito.

\textbf{C) Distractor. }La prueba ya fue anunciada con duración y operación.

\textbf{Repaso:} \hyperref[sec:l3]{Comprensión lectora: participación social}. \textbf{Fuentes:} \cite{ceneval,reading}. \hyperref[q:L-029]{Volver al reactivo}.\par\medskip

\subsection*{L-030 \quad Respuesta A}\label{s:L-030}
\textbf{Razonamiento.} El cierre pide incluir costo, cifra ausente del comunicado inicial.

\textbf{B) Distractor. }No se pide identificar a todos los habitantes.

\textbf{C) Distractor. }No guarda relación con la nota de residuos.

\textbf{Repaso:} \hyperref[sec:l3]{Comprensión lectora: participación social}. \textbf{Fuentes:} \cite{ceneval,reading}. \hyperref[q:L-030]{Volver al reactivo}.\par\medskip

\section{Redacción indirecta: ámbito de estudio}

\subsection*{E-001 \quad Respuesta B}\label{s:E-001}
\textbf{Razonamiento.} Identifica comparación, variable, participantes y contexto de tarea.

\textbf{A) Distractor. }No define acción ni resultado específico.

\textbf{C) Distractor. }Es una afirmación absoluta sin delimitación comprobable.

\textbf{Repaso:} \hyperref[sec:e1]{Redacción indirecta: ámbito de estudio}. \textbf{Fuentes:} \cite{rae,writing}. \hyperref[q:E-001]{Volver al reactivo}.\par\medskip

\subsection*{E-002 \quad Respuesta C}\label{s:E-002}
\textbf{Razonamiento.} El sujeto coordinado exige plural y el atributo concuerda en masculino plural con géneros distintos.

\textbf{A) Distractor. }Usa singular para un sujeto plural.

\textbf{B) Distractor. }El participio singular femenino no concuerda con el conjunto.

\textbf{Repaso:} \hyperref[sec:e1]{Redacción indirecta: ámbito de estudio}. \textbf{Fuentes:} \cite{rae,writing}. \hyperref[q:E-002]{Volver al reactivo}.\par\medskip

\subsection*{E-003 \quad Respuesta B}\label{s:E-003}
\textbf{Razonamiento.} Confirmar introduce aquí una subordinada con que, sin preposición.

\textbf{A) Distractor. }La preposición sola no introduce correctamente esa subordinada.

\textbf{C) Distractor. }Agrega una preposición no exigida: dequeísmo.

\textbf{Repaso:} \hyperref[sec:e1]{Redacción indirecta: ámbito de estudio}. \textbf{Fuentes:} \cite{rae,writing}. \hyperref[q:E-003]{Volver al reactivo}.\par\medskip

\subsection*{E-004 \quad Respuesta C}\label{s:E-004}
\textbf{Razonamiento.} Haber con valor existencial es impersonal y se usa en singular.

\textbf{A) Distractor. }Pluraliza indebidamente el uso existencial.

\textbf{B) Distractor. }También pluraliza ese uso impersonal.

\textbf{Repaso:} \hyperref[sec:e1]{Redacción indirecta: ámbito de estudio}. \textbf{Fuentes:} \cite{rae,writing}. \hyperref[q:E-004]{Volver al reactivo}.\par\medskip

\subsection*{E-005 \quad Respuesta A}\label{s:E-005}
\textbf{Razonamiento.} Cómo lleva tilde en interrogativa indirecta; el artículo el no la lleva.

\textbf{B) Distractor. }Omite la tilde del interrogativo indirecto.

\textbf{C) Distractor. }Él con tilde es pronombre, no artículo ante informe.

\textbf{Repaso:} \hyperref[sec:e1]{Redacción indirecta: ámbito de estudio}. \textbf{Fuentes:} \cite{rae,writing}. \hyperref[q:E-005]{Volver al reactivo}.\par\medskip

\subsection*{E-006 \quad Respuesta C}\label{s:E-006}
\textbf{Razonamiento.} La repetición explícita identifica de quién es la propuesta.

\textbf{A) Distractor. }Su puede referirse a cualquiera de los dos participantes.

\textbf{B) Distractor. }La formulación no resuelve de forma inequívoca el primer posesivo.

\textbf{Repaso:} \hyperref[sec:e1]{Redacción indirecta: ámbito de estudio}. \textbf{Fuentes:} \cite{rae,writing}. \hyperref[q:E-006]{Volver al reactivo}.\par\medskip

\subsection*{E-007 \quad Respuesta A}\label{s:E-007}
\textbf{Razonamiento.} Los dos puntos anuncian la lista y las comas separan sus miembros.

\textbf{B) Distractor. }Separa inadecuadamente el verbo de su complemento y no delimita la lista.

\textbf{C) Distractor. }No separa los primeros elementos y añade una coma injustificada.

\textbf{Repaso:} \hyperref[sec:e1]{Redacción indirecta: ámbito de estudio}. \textbf{Fuentes:} \cite{rae,writing}. \hyperref[q:E-007]{Volver al reactivo}.\par\medskip

\subsection*{E-008 \quad Respuesta B}\label{s:E-008}
\textbf{Razonamiento.} Sin embargo contrapone la limitación y la utilidad observada.

\textbf{A) Distractor. }Introduce explicación o reformulación.

\textbf{C) Distractor. }Introduce consecuencia, no el contraste solicitado.

\textbf{Repaso:} \hyperref[sec:e1]{Redacción indirecta: ámbito de estudio}. \textbf{Fuentes:} \cite{rae,writing}. \hyperref[q:E-008]{Volver al reactivo}.\par\medskip

\subsection*{E-009 \quad Respuesta A}\label{s:E-009}
\textbf{Razonamiento.} La formulación limita el alcance y no confunde asociación con certeza universal.

\textbf{B) Distractor. }Excede la evidencia de una muestra limitada.

\textbf{C) Distractor. }La metodología sí importa para interpretar resultados.

\textbf{Repaso:} \hyperref[sec:e1]{Redacción indirecta: ámbito de estudio}. \textbf{Fuentes:} \cite{rae,writing}. \hyperref[q:E-009]{Volver al reactivo}.\par\medskip

\subsection*{E-010 \quad Respuesta B}\label{s:E-010}
\textbf{Razonamiento.} Primero se plantea la incertidumbre, luego el método y finalmente se anuncian resultados.

\textbf{A) Distractor. }Por ello requiere un antecedente que aún no aparece.

\textbf{C) Distractor. }Anuncia resultados antes de establecer problema y acción.

\textbf{Repaso:} \hyperref[sec:e1]{Redacción indirecta: ámbito de estudio}. \textbf{Fuentes:} \cite{rae,writing}. \hyperref[q:E-010]{Volver al reactivo}.\par\medskip

\subsection*{E-011 \quad Respuesta C}\label{s:E-011}
\textbf{Razonamiento.} No se intercala una coma entre el sujeto completo y su verbo.

\textbf{A) Distractor. }Separa sin justificación sujeto y predicado.

\textbf{B) Distractor. }Abre un inciso sin cerrarlo y altera la estructura.

\textbf{Repaso:} \hyperref[sec:e1]{Redacción indirecta: ámbito de estudio}. \textbf{Fuentes:} \cite{rae,writing}. \hyperref[q:E-011]{Volver al reactivo}.\par\medskip

\subsection*{E-012 \quad Respuesta B}\label{s:E-012}
\textbf{Razonamiento.} Sino contrapone términos; si no introduce una condición negativa.

\textbf{A) Distractor. }La condición requiere si no separado.

\textbf{C) Distractor. }Intercambia las dos construcciones.

\textbf{Repaso:} \hyperref[sec:e1]{Redacción indirecta: ámbito de estudio}. \textbf{Fuentes:} \cite{rae,writing}. \hyperref[q:E-012]{Volver al reactivo}.\par\medskip

\subsection*{E-013 \quad Respuesta C}\label{s:E-013}
\textbf{Razonamiento.} Separa observación e interpretación y marca su incertidumbre.

\textbf{A) Distractor. }Afirma causalidad sin evidencia adicional.

\textbf{B) Distractor. }La preferencia no constituye evidencia causal.

\textbf{Repaso:} \hyperref[sec:e1]{Redacción indirecta: ámbito de estudio}. \textbf{Fuentes:} \cite{rae,writing}. \hyperref[q:E-013]{Volver al reactivo}.\par\medskip

\subsection*{E-014 \quad Respuesta A}\label{s:E-014}
\textbf{Razonamiento.} Aún equivale a todavía; aun con equivale aquí a incluso con.

\textbf{B) Distractor. }Intercambia la tilde en ambos usos.

\textbf{C) Distractor. }En la concesión equivalente a incluso no corresponde aún.

\textbf{Repaso:} \hyperref[sec:e1]{Redacción indirecta: ámbito de estudio}. \textbf{Fuentes:} \cite{rae,writing}. \hyperref[q:E-014]{Volver al reactivo}.\par\medskip

\subsection*{E-015 \quad Respuesta C}\label{s:E-015}
\textbf{Razonamiento.} Explica la idea con lenguaje accesible y conserva límites sin prometer resultados universales.

\textbf{A) Distractor. }El tecnicismo sin explicación dificulta el propósito divulgativo indicado.

\textbf{B) Distractor. }Atribuye propiedades falsas y absolutas a la herramienta.

\textbf{Repaso:} \hyperref[sec:e1]{Redacción indirecta: ámbito de estudio}. \textbf{Fuentes:} \cite{rae,writing}. \hyperref[q:E-015]{Volver al reactivo}.\par\medskip

\section{Redacción indirecta: participación social}

\subsection*{E-016 \quad Respuesta B}\label{s:E-016}
\textbf{Razonamiento.} Declara una solicitud respetuosa y anuncia la justificación pertinente.

\textbf{A) Distractor. }Atribuye una decisión que no se ha comunicado.

\textbf{C) Distractor. }No adecua registro ni fundamenta la petición.

\textbf{Repaso:} \hyperref[sec:e2]{Redacción indirecta: participación social}. \textbf{Fuentes:} \cite{rae,writing}. \hyperref[q:E-016]{Volver al reactivo}.\par\medskip

\subsection*{E-017 \quad Respuesta C}\label{s:E-017}
\textbf{Razonamiento.} El sujeto plural requiere pueden; interesados concuerda con el conjunto.

\textbf{A) Distractor. }El adjetivo y el verbo están indebidamente en singular.

\textbf{B) Distractor. }El verbo singular no concuerda con el sujeto.

\textbf{Repaso:} \hyperref[sec:e2]{Redacción indirecta: participación social}. \textbf{Fuentes:} \cite{rae,writing}. \hyperref[q:E-017]{Volver al reactivo}.\par\medskip

\subsection*{E-018 \quad Respuesta B}\label{s:E-018}
\textbf{Razonamiento.} La coma separa el nombre de la persona a quien se dirige la petición.

\textbf{A) Distractor. }Separa artículo y sustantivo.

\textbf{C) Distractor. }Separa verbo y complemento.

\textbf{Repaso:} \hyperref[sec:e2]{Redacción indirecta: participación social}. \textbf{Fuentes:} \cite{rae,writing}. \hyperref[q:E-018]{Volver al reactivo}.\par\medskip

\subsection*{E-019 \quad Respuesta C}\label{s:E-019}
\textbf{Razonamiento.} Tú pronombre lleva tilde; tu posesivo no.

\textbf{A) Distractor. }Invierte los usos de pronombre y posesivo.

\textbf{B) Distractor. }El segundo tu es posesivo ante solicitud.

\textbf{Repaso:} \hyperref[sec:e2]{Redacción indirecta: participación social}. \textbf{Fuentes:} \cite{rae,writing}. \hyperref[q:E-019]{Volver al reactivo}.\par\medskip

\subsection*{E-020 \quad Respuesta A}\label{s:E-020}
\textbf{Razonamiento.} Porque introduce la causa de no asistir.

\textbf{B) Distractor. }Por qué corresponde a una interrogación o estructura distinta.

\textbf{C) Distractor. }Porqué es un sustantivo y aquí no tiene determinante ni función nominal.

\textbf{Repaso:} \hyperref[sec:e2]{Redacción indirecta: participación social}. \textbf{Fuentes:} \cite{rae,writing}. \hyperref[q:E-020]{Volver al reactivo}.\par\medskip

\subsection*{E-021 \quad Respuesta C}\label{s:E-021}
\textbf{Razonamiento.} Se identifica a quien emite el juicio sin convertirlo en hecho universal.

\textbf{A) Distractor. }Presenta una generalización sin respaldo ni atribución.

\textbf{B) Distractor. }Formula una valoración absoluta como si fuera definición.

\textbf{Repaso:} \hyperref[sec:e2]{Redacción indirecta: participación social}. \textbf{Fuentes:} \cite{rae,writing}. \hyperref[q:E-021]{Volver al reactivo}.\par\medskip

\subsection*{E-022 \quad Respuesta A}\label{s:E-022}
\textbf{Razonamiento.} Especifica acción, objeto, lugar y plazo.

\textbf{B) Distractor. }Los referentes y el tiempo son imprecisos.

\textbf{C) Distractor. }Depende de información no explicitada.

\textbf{Repaso:} \hyperref[sec:e2]{Redacción indirecta: participación social}. \textbf{Fuentes:} \cite{rae,writing}. \hyperref[q:E-022]{Volver al reactivo}.\par\medskip

\subsection*{E-023 \quad Respuesta B}\label{s:E-023}
\textbf{Razonamiento.} Ha es auxiliar de enviado; a introduce destinatario.

\textbf{A) Distractor. }El segundo ha debe ser la preposición a.

\textbf{C) Distractor. }Intercambia auxiliar y preposición.

\textbf{Repaso:} \hyperref[sec:e2]{Redacción indirecta: participación social}. \textbf{Fuentes:} \cite{rae,writing}. \hyperref[q:E-023]{Volver al reactivo}.\par\medskip

\subsection*{E-024 \quad Respuesta A}\label{s:E-024}
\textbf{Razonamiento.} Indica con precisión el cambio y confirma las condiciones vigentes.

\textbf{B) Distractor. }Sugiere cambios indeterminados y no informa lo necesario.

\textbf{C) Distractor. }No comunica que esa sea la nueva sede.

\textbf{Repaso:} \hyperref[sec:e2]{Redacción indirecta: participación social}. \textbf{Fuentes:} \cite{rae,writing}. \hyperref[q:E-024]{Volver al reactivo}.\par\medskip

\subsection*{E-025 \quad Respuesta B}\label{s:E-025}
\textbf{Razonamiento.} Se presenta la petición, se fundamenta y se cierra.

\textbf{A) Distractor. }Cierra antes de formular la petición principal.

\textbf{C) Distractor. }La despedida no introduce eficazmente el trámite.

\textbf{Repaso:} \hyperref[sec:e2]{Redacción indirecta: participación social}. \textbf{Fuentes:} \cite{rae,writing}. \hyperref[q:E-025]{Volver al reactivo}.\par\medskip

\subsection*{E-026 \quad Respuesta C}\label{s:E-026}
\textbf{Razonamiento.} Ahí indica lugar, hay existencia y ay es interjección.

\textbf{A) Distractor. }Ahí no expresa la exclamación de dolor indicada.

\textbf{B) Distractor. }Ay no introduce el lugar y hay no es la interjección.

\textbf{Repaso:} \hyperref[sec:e2]{Redacción indirecta: participación social}. \textbf{Fuentes:} \cite{rae,writing}. \hyperref[q:E-026]{Volver al reactivo}.\par\medskip

\subsection*{E-027 \quad Respuesta B}\label{s:E-027}
\textbf{Razonamiento.} Define condición, mecanismo y fecha de selección.

\textbf{A) Distractor. }No permite conocer cómo se aplicará la selección.

\textbf{C) Distractor. }No establece procedimiento.

\textbf{Repaso:} \hyperref[sec:e2]{Redacción indirecta: participación social}. \textbf{Fuentes:} \cite{rae,writing}. \hyperref[q:E-027]{Volver al reactivo}.\par\medskip

\subsection*{E-028 \quad Respuesta C}\label{s:E-028}
\textbf{Razonamiento.} Las tres acciones están expresadas como infinitivos coordinados.

\textbf{A) Distractor. }Mezcla formas sin una estructura uniforme.

\textbf{B) Distractor. }La combinación produce una lista mal construida.

\textbf{Repaso:} \hyperref[sec:e2]{Redacción indirecta: participación social}. \textbf{Fuentes:} \cite{rae,writing}. \hyperref[q:E-028]{Volver al reactivo}.\par\medskip

\subsection*{E-029 \quad Respuesta A}\label{s:E-029}
\textbf{Razonamiento.} Conserva el dato explícito sin inferir cuántas solicitudes había ni qué proporción fue atendida.

\textbf{B) Distractor. }Contradice las solicitudes pendientes.

\textbf{C) Distractor. }Niega sin respaldo las atenciones efectivas.

\textbf{Repaso:} \hyperref[sec:e2]{Redacción indirecta: participación social}. \textbf{Fuentes:} \cite{rae,writing}. \hyperref[q:E-029]{Volver al reactivo}.\par\medskip

\subsection*{E-030 \quad Respuesta C}\label{s:E-030}
\textbf{Razonamiento.} Formula una posición institucional y aporta una razón que la sustenta.

\textbf{A) Distractor. }Presenta un dato informativo sin desarrollar una postura editorial.

\textbf{B) Distractor. }Corresponde a instrucciones de convocatoria, no a un argumento editorial.

\textbf{Repaso:} \hyperref[sec:e2]{Redacción indirecta: participación social}. \textbf{Fuentes:} \cite{rae,writing}. \hyperref[q:E-030]{Volver al reactivo}.\par\medskip